%\docume  ntclass[11pt,tightenlines,nofootinbib,longbibliography]{revtex4-2}
%\documentclass[11pt,reprint,prl,twocolumn]{revtex4-2}
\documentclass[11pt,letterpaper]{article}
%\usepackage[T1]{fontenc}
\usepackage{jheppub}
 
\usepackage{xcolor}
\usepackage{mathdots}
\usepackage{amsmath}
\usepackage{comment}
\usepackage{mathtools}
\usepackage{amssymb}
\usepackage{amsfonts}
\usepackage{amsthm}
\usepackage{mathrsfs}
\usepackage{footnote}
\usepackage{graphicx} 
\usepackage{subcaption}
\usepackage{stmaryrd}
\usepackage[colorlinks=true
,urlcolor=blue
,anchorcolor=blue
,citecolor=blue
,filecolor=blue
,linkcolor=blue
,menucolor=blue
,pagecolor=blue
,linktocpage=true
,pdfproducer=medialab
,pdfa=true
]{hyperref}

\newtheorem{thm}{Theorem}
\newtheorem{lem}{Lemma}
\newtheorem{conj}{Conjecture}
\newtheorem{cor}{Corollary}
\newtheorem{defn}{Definition}
\newtheorem{prop}{Proposition}
\newtheorem{cl}{Claim}
\newtheorem{qn}{Question}
\newtheorem{pr}{Proposition}
\newtheorem{rem}{Remark}


\newcommand{\dd}{\mathrm{d}}
\newcommand{\ket}[1]{| #1 \rangle}
\newcommand{\kket}[1]{| #1 \rangle\rangle}
\newcommand{\bra}[1]{\langle #1 |}
\newcommand{\braket}[2]{\left\langle #1 \mid #2 \right\rangle}
\newcommand{\ketbra}[2]{\left|#1\right\rangle\!\!\left\langle #2\right|}
\def\id{\mathbf{1}}
\def\N{\mathcal{N}}
\def\eps{\varepsilon}
\def\O{\mathcal{O}}
\def\h{\mathcal{H}}
\def\k{\mathcal{K}}
\def\id{\mathbb{I}}
\newcommand*{\Tr}{\mathrm{Tr}}
\def\tr{\mathrm{tr}}
\def\D{\mathcal{D}}
\def\E{\mathcal{E}}
\def\A{\mathcal{A}}
\def\C{\mathcal{C}}
\def\map{\mathcal}
\def\set{\mathsf}

\newcommand{\JW}[1]{\textcolor{teal}{[JW:\ #1]}}

\newcommand{\YY}[1]{\textcolor{blue}{[YY:\ #1]}}




\title{Free entropy and quantum minimum description length}
%\author{ \\
%{\small \it Stanford Institute for Theoretical Physics, Stanford University, Stanford, CA 94305}}
\abstract{
In this work, we explore the operational meaning of free entropy and free entropy dimension in the context of quantum information theory. We find that free entropy (dimension) characterizes how well one can universally describe various quantum objects using quantum states. Explicitly, we discover that for quantum states, free entropy characterizes the optimal state compression without purification; for unitary quantum evolutions, it characterizes the optimal program size for universal unitary programming; and for quantum observables, it characterizes the optimal program size for universal observable programming. We provide convincing evidence in support of the claim that free entropy measures quantum minimum description length.}


\begin{document}

\maketitle


\section{Introduction}
%Free entropy
The free entropy $\chi$ was invented as (another) non-commutative analog to the Shannon entropy~\cite{voiculescu1993analogues,voiculescu1994analogues,voiculescu1996analogues,voiculescu1998analogues,voiculescu1999analogues}.  So far it has only been used as a tool to prove various important results in Type II von Neumann algebra~\cite{voiculescu1996analogues,ge1997applications,ge1998applications}, which is also the primary purpose of its invention. While termed as an ``entropy'', its operational meaning in physics has not yet been found. It is hence the goal of this work to find one good physical interpretation of the free entropy.

Mathematically, we can make a parallel between the free entropy and the von Neumann entropy. Recall that Shannon entropy admits two formulations. The canonical definition is $H(\mathbf{p})=\sum_i -p_i\log p_i$, and the von Neumann entropy is a direct generalization of this formula to the non-commutative setting.

By virtue of Shannon's source coding theorem, there is another formula for the Shannon entropy via the asymptotic equipartition property. Consider a random variable $X$ that takes values from an finite alphabet $\mathcal{K}$ and its probability distribution is given by $\vec p$.
\begin{equation}\label{eq:aep}
    H(\vec{p})=\inf_{\eps>0}\lim_{N\to\infty}\frac1N\log \#\{ x^N\in \mathcal{K}^N \, \big\vert\, ||P_{x^N}-\vec{p}||_1\le\eps \}\ ,
\end{equation}
where $x^N$ is a string of length $N$, $P_{x^N}$ denotes the empirical distribution based on $x^N$ and $\#$ denotes the counting measure. This is the definition that makes the Shannon entropy into the Boltzmann form $H=\log W$. It is the log-counting of the number of typical strings output by a source characterized by $\vec{p}$. 

Free entropy is a non-commutative generalization of \eqref{eq:aep}~\cite{voiculescu2002free}, in which random bit strings are replaced by random matrices. Informally, $\chi(\rho)$ for a bounded self-adjoint random variable $\rho\in(\A,\tau)$ is defined as
\begin{equation}\label{eq:free_ent}
    \chi(\rho):= \inf_{m\in\mathbb{N}}\inf_{\eps>0}\lim_{N\to\infty}\left(\frac{1}{N^2}\log \mathrm{Vol}\{ X\in M^{\mathrm{s.a.}}_N(\mathbb{C})\, \big\vert\, |\tr X^m-\tau(\rho^m)|\le\eps \}+\frac12\log N\right)\ ,
\end{equation}
where $M^{\mathrm{s.a.}}_N(\mathbb{C})$ denotes the set of $N\times N$ self-adjoint matrices over complex numbers, and $\tr:=\tr/N$ denotes the normalized trace. The volume measure can be chosen to be the Lebesgue measure over $\mathbb{R}^{N^2}$,\footnote{One can also choose the Gaussian measure, and the resulting formula is identical up to some additive constant.} in which $M^{\mathrm{s.a.}}_N(\mathbb{C})$ can be embedded. The $\frac12\log N$ term removes an universal divergence in the ${N^{-2}}\log \mathrm{Vol}$ term. It can be shown that~\cite{voiculescu1994analogues,hiai2000semicircle},
\begin{equation}\label{eq:free_ent_formula}
    \chi(\rho) = \int_{\mathbb{R}}\int_{\mathbb{R}} \log|x-y|\mu_{\rho}(x)\mu_{\rho}(y) \dd x \dd y + \frac34+\frac12\log 2\pi\ ,
\end{equation}
where $\mu_{\rho}$ denotes the spectral density of $\rho$, which is determined by the moments $\tau(\rho^m)$. Up to the constant terms that are independent of the spectrum, the logarithmic term comes from the logarithm of the Vandermonde determinant. 


Free entropy theory is mostly formulated in infinite dimensions for random variables in a Type II von Neumann algebra. In finite dimensions, $\rho\in (M_d(\mathbb{C}),\tr)$, and $\mu_\rho(x) = \frac1d\sum_{i=1}^d\delta(x-p_i)$ is a sum of atomic measures at the eigenvalues $p_i$ of $\rho$. We have
\begin{equation}
    \chi(\rho) = \frac{2}{d^2}\sum_{i< j} \log|p_i-p_j|-\infty\ .
\end{equation}
(More details on this below.) The divergence due to $\sum_i g_i$ collisions can be renormalized out by considering the free entropy dimension. The resulting finite piece only sums over distinct pairs of eigenvalues. We regard the finite piece as the appropriate formula of the free entropy in finite dimensions. In general, when the spectrum of $\rho$ is degenerate, we have
\begin{equation}\label{eq:free_ent_formula_finite}
    \chi_\mathrm{reg}(\rho) = \frac{2}{d^2}\sum_{i< j} g_ig_j\log|p_i-p_j|\ ,
\end{equation}
where $i$ and $j$ label distinct eigenvalues of $\rho$.\\


The definition also applies to unitaries, and the formula is identical up to changing the integration contour to the unit circle $\mathbb{T}$~\cite{hiai2000semicircle}
\begin{equation}\label{eq:free_ent_formula_unitary}
    \chi(u) = \int_{\mathbb{T}}\int_{\mathbb{T}} \log|z-z'| \mu_u(z) \mu_u(z') \dd z \dd z'\ .
\end{equation}
This quantity could potentially be relevant in characterizing the optimal memory size for unitary programming~\cite{yang2020optimal}.

The multivariate free entropy is also defined in the same way as \eqref{eq:free_ent}, where one demands the moments of all the words to be matched by some matrix models at large $N$. In the multivariate case, explicit general formulas are not available. However, a formula is known for the case of two projections~\cite{hiai2005large,hiai2006notes,hiai2006log}.

Entropy is meant to be an extensive quantity in physics. The von Neumann entropy is extensive in the sense that it is additive under tensor product, which is the appropriate classical (commutative) notion of how two unrelated systems should be joined. In contrast, the free entropy is additive under the free product instead.\footnote{Interestingly, free independence and tensor independence are the only two viable notions of independence in non-commutative probability theory under some generic axioms~\cite{speicher1996universal,muraki2002five}. This fact makes the parallel between von Neumann entropy and free entropy more natural.} \\

\section{Defining free entropy in finite dimensions}
According to the mathematical definition~\eqref{eq:free_ent}, the free entropy takes the value of $-\infty$ whenever the spectrum has an atom. It is mostly interesting for random variables with a continuous spectrum.\footnote{In this sense, it is more analogous to the differential Shannon entropy than the (discrete) Shannon entropy.} For our purposes, however, it is crucial to adjust the standard definition for finite-dimensional density operators, such that one can compare free entropies of different states. 

To this end, we propose two principled approaches to obtain \eqref{eq:free_ent_formula_finite}. The first approach follows how physicists usually deal with infinities. The second approach is more straightforward, but it involves altering the standard definition~\eqref{eq:free_ent}.

\subsection{Free entropy as a renormalized entropy}
Let's recall Shannon's differential entropy of a random variable $X$ with distribution density function $f(x)$,
\begin{equation}
    h(X) := \int -f(x)\log f(x)\dd x.
\end{equation}

While it is a mathematically sound definition, as compared to the Shannon entropy, it suffers several flaws as a physical quantity: It can become negative, it is not invariant under a change of variables and it is not dimensionally correct. 

It's well understood that the differential entropy should be viewed as a \emph{renormalized entropy}. Suppose that $X$ takes values in $\mathbb{R}$, one can discretize the support of $f(x)$ into bins of width $\eps$. If $f(x)\log f(x)$ is Riemann integrable, one can show that (cf. e.g. Chapter 9 in Thomas and Cover.)
\begin{equation}
   \lim_{\eps\to 0} H(X;\eps)+\log\eps = h(X)
\end{equation}
where $H(X;\eps)$ is the Shannon entropy of the discretized random variable. It is divergent as $\log\eps^{-1}$, so one removes this divergence by introducing a counter-term $\log\eps$ and arrives at the finite value $h(X)$. In other words, the physically meaningful entropy Shannon entropy contains a universal divergent piece and a finite piece that depends on the distribution density. 
\begin{equation}
    H(X;\eps)\simeq\log\eps^{-1} + h(X)
\end{equation}

In this regard, the free entropy is not too different from the differential entropy. It is itself a renormalized entropy. When looking for a physical meaning, one should restore the counting measure. We can write Vol$\approx\eps^{N^2}\#$ in~\eqref{eq:free_ent}, where $\eps^{N^2}$ denotes the unit volume of a cell of length $\eps$ in Lebesgue measure. The exact geometric details only shift the entropy by constants and are hence irrelevant for the entropy measure. Let's use $\chi(X;\eps)$ to denote free entropy measured by counting the cells in the $\eps$-net of $M_N^\mathrm{s.a.}(\mathbb{C})$, we have
\begin{equation}\label{eq:physical_free_ent}
    \chi_\mathrm{phy}(\rho;\eps)\simeq\log\eps^{-1} + \chi(\rho)
\end{equation}
We claim that $\chi_\mathrm{phy}(\rho;\eps)$ is the \emph{physical free entropy}, where $\eps$ shall be chosen according to the relevant precision considered in the operational task.\\

Let's now consider $\chi(\rho)$ for a density operator $\rho$ with spectral density
\begin{equation}
    \mu(x):= \frac1d\sum_i\delta(x-p_i)\ .
\end{equation}
To regularize the infinity encountered in the collisions in~\eqref{eq:free_ent_formula}, we deform the integration contour by adding to $\rho$ a semicircular of $\eps$ width.  That is we compute the free entropy $\chi(\rho+\eps s)$, which yields a finite value. We set $s$ free from $\rho$. $\eps$ is chosen so that the distance between $\rho+\eps s$ and $\rho$ is of the order $\eps$. 

The spectral density of $\rho+\eps s$ can be calculated using the free additive convolution. Its support has as many disconnected domains as the number of distinct eigenvalues, and each domain has a width of order $\eps$. The resulting free entropy admits the following expansion at small $\eps$
\begin{equation}
    \chi(\rho+\eps s) \simeq \frac{\sum_ig_i^2}{d^2}\log\eps + \frac{2}{d^2}\sum_{i<j}g_ig_j\log |p_i-p_j| + \O(\eps)
\end{equation}
where $i,j$ labels distinct eigenvalues of $\rho$. We define the second term as $\chi_\mathrm{reg}(\rho)$ according to~\eqref{eq:free_ent_formula_finite}.

This regularized free entropy identifies a divergent piece in $\chi(\rho)$. We can now renormalize it in the physical entropy by absorbing the divergent piece into the first term in~\eqref{eq:physical_free_ent}. 
\begin{equation}
    \chi_\mathrm{phy}(\rho;\eps)\simeq\left(1-\frac{\sum_ig_i^2}{d^2}\right)\log\eps^{-1} + \chi_\mathrm{reg}(\rho) + \O(\eps).
\end{equation}
Hence, the physical free entropy $\chi_\mathrm{phy}(\rho;\eps)$ is a renormalized quantity. The bare divergence in $\chi(\rho)$ can be absorbed into the universal divergent term.

The coefficient of the leading divergent piece is called \emph{free entropy dimension} of $\rho$, which characterizes the growth of $\chi_\mathrm{phy}(\rho;\eps)$ with $\eps\to 0$. (cf. Sec.\ref{sec:fed} below) \\

What is the maximal free entropy and what is the maximizing eigenvalue distribution? In finite dimensions, it is clear that a non-degenerate spectrum is favored. It remains to find the maximal set of $\{p_i\}$. Consider the following multivariate function
\begin{equation}
    f(p_1,\ldots,p_{d-1},\alpha) = 2\sum_{i<j}\log |p_i-p_j| + \alpha (1-\sum_{i=1}^d p_i)
\end{equation}
We let $p_i>p_j$ for $i<j$ and set $p_d=0$. The saddle point amounts to solving the following system of polynomial equations.
\begin{equation}
   \forall k, 2\sum_{i, i\neq k}(p_k-p_i)^{-1} = \alpha
\end{equation}
under the constraint $\sum_{i=1}^d p_i=1$. The solution to this optimization problem is called the Fekete points. In the literature, it's common to find the Fekete points for some given interval, say $[-1,1]$ or $[0,1]$. What we want are the Fekete points on a simplex. 

At large $d$,  the profile of the Fekete points should be well approximated by a continuous density function. The solution to this maximization problem has essentially been solved by Petz and Hiai~\cite{hiai1998maximizing}. We only need to adapt their solution to our setting of density operators.

\subsection{A modified definition in finite dimensions}
The motivation of the large $N$ limit in~\eqref{eq:free_ent} is that one would like to better approximate the given (continuous) spectrum with large matrices. When the spectrum is finitely discrete like for a $d$-dimensional density operator, however, matrices of size $N=d$ are sufficient to exactly reproduce the spectrum of the given density operator. We can then omit the large $N$ limit in~\eqref{eq:free_ent} and directly set $N=d$. Also, we can demand the error to be zero.

It's then natural to define the finite-dimensional free entropy simply as the log-volume of the unitary orbit. ...


\subsection{Free entropy dimension}\label{sec:fed}
When free entropy diverges, the submanifold of matricial microstates has Lebesgue measure zero. This means that the submanifold has fewer dimensions than $N^2$. In this case, it is important to first describe the dimension of the space spanned by the matricial microstates before we discuss their appropriate volumes. This motivates the definition of \emph{free entropy dimension} that generalizes the Minkowski dimension (or entropy dimension) to the non-commutative setting. 

The free entropy dimension is defined as 
\begin{equation}\label{eq:fed}
    \delta(a_1,\ldots,a_n):=n+\lim_{\eps\to 0}\frac{\chi(a_1+\eps S_1,\ldots,a_n+\eps S_n)}{-\log\eps}
\end{equation}
where $\{a_1,\ldots,a_n\}$, and the semicirculars $ S_1,\ldots, S_n$ are free. The semicircular smearing helps to remove any logarithmic divergences from the free entropy, and the limit extracts the degree of the divergences and subtract them from $n$.
We have $0\leq\delta(a_1,\dots,a_n)\leq n$ and $\delta(a_1,\ldots,a_n)=n$ iff $\chi(a_1,\ldots,a_n)>-\infty$. 

Note that the free entropy dimension is exactly the factor controlling the leading $\log\epsilon^{-1}$ term in the definition of physical free entropy.
\begin{equation}
    \chi_\mathrm{phy}(a;\epsilon)=\delta(a)\log\epsilon^{-1} + \chi_\mathrm{reg}(a) + \mathcal{O}(\epsilon)
\end{equation}
where $\chi_\mathrm{reg}$ is the residual finite part of the free entropy. 

In light of this, we also can define the physical free entropy for multiple variables as
\begin{equation}
    \chi_\mathrm{phy}(a_1,\ldots,a_n;\epsilon):=\delta(a_1,\ldots,a_n)\log\epsilon^{-1} + \chi_\mathrm{reg}(a_1,\ldots,a_n)\ .
\end{equation}


For a single variable, the free entropy dimension has a formula that substracts from $1$ the contributions from the atoms
\begin{equation}
    \delta(a)=1-\sum_{t\in\mathbb{R}} \mu_a(\{t\})^2
\end{equation}
where $\mu_a$ is the distribution of $a\in (\mathcal{A},\tau)$. Using this formula, we can work out some simple examples: For a projection $P\in (\mathcal{A},\tau)$, we have
\begin{equation}
    \delta(P)=2\tau(P)[1-\tau(P)]
\end{equation}
For a rank-$k$ $d$-dimensional projector $P\in (\mathcal{M}_d,\tr)$, we have
\begin{equation}
    \delta(P)=2\frac{k}d\left[1-\frac{k}d\right]\ .
\end{equation}
For a $d$-dimensional Hermitian operator $H$ that has degeneracies $g_i$, where $i$ lables distinct eigenvalues, we have
\begin{equation}
    \delta(H)=1 - \sum_i\frac{g_i^2}{d^2}\ .
\end{equation}


\section{Free entropy of density operators: Quantum state compression}

We are looking for a compression task in quantum theory, such that the minimal memory size is characterized by the free entropy. Since free entropy measures the number of matrices that match a given spectrum, we find the following scenario studied by Yang-Chiribella-Hayashi (YCH) the most promising setup~\cite{yang2016optimal}.

The task is to compress a source that outputs i.i.d. copies of a density matrix $\rho$, $\rho^{\otimes n}$. We assume that the spectrum of $\rho$ is known but the eigenvectors are not. Compared to Schumacher compression, the error is judged without including the purification. Let the memory be~$M$. We say $(\E,\D)$ is a $(|M|,\delta)$-code if
\begin{equation}\label{eq:error}
    \frac12||\rho^{\otimes n}-\D\circ\E(\rho^{\otimes n})||_1\le\delta\ .
\end{equation} 

The goal is to determine the optimal code with the minimal rate $\log M$ with a vanishing $\delta$ at large $n$. We would like to compute the asymptotic rate as an expansion in $n$ up to the terms that vanish at large $n$. That is, we want to determine both the leading $\O(\log n)$ term and the finite $\O(1)$ term in the expansion. Although the leading rate is found by YCH for qubits~\cite{yang2016optimal} and later by YBCH for qudits~\cite{yang2018compression}, the spectrum-dependent $\O(1)$ correction is still open. 

Our conjecture is that the minimal memory size is determined by the free entropy~\eqref{eq:free_ent_formula_finite} up to universal terms. That is, when we compare the memory sizes for different density matrix spectra, the difference is determined by the free entropy difference.

We give some heuristics for our conjecture. Consider an $\eps$-net $\Gamma$ that discretizes the unitary orbit of a given density matrix spectrum $\mathbf{p}$ that is non-degenerate. The volume is roughly the number of cells $|\Gamma|$ in the net times the unit cell volume
\begin{equation}
    \mathrm{Vol}(\rho)=C |\Gamma| \eps^{d(d-1)}\ ,
\end{equation}
where $C$ is some measure-dependent geometric constant and $d(d-1)$ is the number of degrees of freedom (dimensions) of the orbit. In general, with degeneracies $g_i$, we have $d(d-1)-\sum_i g_i(g_i-1)$ dimensions. On the other hand, we have the free entropy~\eqref{eq:free_ent_formula_finite} as the log-volume,
\begin{equation}
    \mathrm{Vol}(\rho) \approx c\, 2^{d^2\chi(\rho)}
\end{equation}
with some constant $c$. It follows that 
\begin{equation}
    \log|\Gamma| \simeq d(d-1)\log\eps^{-1} +d^2\chi(\rho) - \log C+\log c\ .
\end{equation}
For a degenerate spectrum, the leading factor changes to $d(d-1)-\sum_i g_i(g_i-1)=d^2-\sum_i g_i^2$.
%Since $\eps$ sets the resolution for $\rho$, then it only makes sense to demand \eqref{eq:error} no finer than a $\delta\simeq \eps$ precision. We will later show that the optimal protocol yields an error in~\eqref{eq:error} that scales inversely with $n$. It follows that 
%\begin{equation}
%    \log|\Gamma| \simeq d(d-1)\log n +d^2\chi(\rho) - \log C+\log c\ .
%\end{equation}


This is the number of bits needed to label different states on the unitary orbit given the spectrum, and it should characterize the minimal memory size. In fact, it turns out a quantum memory can do \emph{quadratically} better in terms of the dimension, and we will show the quantum memory size is 
\begin{equation}
    \log |M| \simeq \frac12\left[(d^2-\sum_i g_i^2))\log n +d^2\chi_\mathrm{reg}(\rho)\right] + \mathrm{Const}=\frac{d^2}{2}\chi_\mathrm{phy}(\rho;n^{-1}) + \mathrm{Const}.
\end{equation}
The factor of $d^2$ here is due to the normalization convention.

\JW{Why does the error scale like $\eps\sim n^{-1}$?}


%\subsection{On covariance of the compression protocol}
%Consider any compression protocol with the encoder-decoder pair $(\mathcal{E},\mathcal{D})$. Denote by $\epsilon(\mathcal{E},\mathcal{D})$ the error of this protocol.


\subsection{Review of YCH on qubits compression}
The YCH scheme for qubits already achieves what we expect. Quoted directly from YCH, their scheme is as follows. 

``The encoder: Perform the Schur transform $V: \h^{\otimes n}\to\bigoplus_{\lambda} V_\lambda\otimes\mathcal{V}_\lambda$,
\begin{equation}\label{eq:schur}
    V(U(g)\rho U(g)^\dagger)^{\otimes n} V^\dagger = \bigoplus_{J=0}^{n/2} q_J \rho_{g,J}\otimes \tau_{J}
\end{equation}
where $g\in\mathrm{U}(2)$, $U(g)$ is the fundamental irrep, $\rho_{g,J}:=U_J(g) \rho_J U_J\dagger(g)$ and $U_J(g)$ is the $2J+1$ dimensional-irrep of $\mathrm{U}(2)$, and $\tau_J$ is the maximally mixed state on the $\mathcal{V}_J$ irrep of $S_n$ that has dimension $(2J+1)\binom{n+1}{n/2+J+1}/(n+1)$. The distribution $q_J$ reads
\begin{equation}
    q_J := \frac{2J+1}{2J^*}\left[B(n/2+J+1;p,n)-B(n/2-J;p,n)\right], B(k;p,n):=\binom{n}{k}p^{k}(1-p)^{n-k}\ .
\end{equation}

The Young diagrams $(n/2+J,n/2-J)$ are labeled by $J$. Then, measure the quantum number~$J$ with the non-demolition measurement that preserves the quantum information in each subspace $\mathcal{U}_J\otimes\mathcal{V}_J$. Let $J^*$ be the most typical outcome, i.e. that maximizes $q_J$. Discard the multiplicity register and apply the cloning channel $\C_{J\to J^*}$ to the remaining state $\rho_{g,J}$. Store the output state into a quantum memory of dimension $2J^* + 1$.

The decoder: Pick a value $K$ at random with probability $q_K$ and apply the cloning channel $\C_{J^*\to K}$ to the
quantum memory. Append a multiplicity register in the maximally mixed state $\tau_K$. Finally, perform the inverse of the Schur transform.''

YCH showed that their scheme has an asymptotically vanishing error $\eps\sim\O(n^{-\frac12})$.

A key ingredient in their scheme is the cloning channel~\cite{werner1998optimal,keyl1999optimal} (extended to include appropriate contractions). The idea is to view $\mathcal{U}_J$ as the symmetric subspace of some $2J$ qubits. It's defined as 
\begin{equation}\label{eq:qubitcloning}
\C_{J\to K}(\rho_J):=
    \begin{cases}
       \frac{2J+1}{2K+1} P_K(\rho_J\otimes I_{2(K-J)} )P_K,\quad J\le K\\
       \tr_{2(J-K)}\rho_J,\qquad\qquad\qquad\quad J>K
    \end{cases}
\end{equation}
mapping the symmetric subspace of $2J$ qubits to the symmetric subspace of $2K$ qubits.

Here the cloning channel is used to map any output on the $J$-th block to the most typical $J^*$ block. One upshot is that it is U(2)-covariant.



We have $J^*=(p-1/2)(n+1)$, so the achievable rate at large $n$ is
\begin{equation}\label{eq:achievable_qubit}
    \log (2J^*+1) \approx \log n + \log (2p-1)\ .
\end{equation}

\subsection{A Generalized Cloning map}

We want to generalize the YCH qubit scheme to qudits. Finding an appropriate U(d)-covariant generalization of the cloning channel that maps any measured Young diagram $\lambda$ to the most typical Young diagram $\lambda^*$, ($\lambda^*_i/n\to p_i$).

We say a quantum channel $\N$ is $UV$-covariant (equivariant) if 
\begin{equation}\label{eq:covariance}
    \N(U\rho U^\dagger) = V\N(\rho)V^\dagger, \quad\forall\rho\ ,
\end{equation}
where the unitaries $U$ and $V$ are representations of U(d) and they act on the input and output Hilbert spaces respectively. We are mostly interested when they act irreducibly on the irreps, so we henceforth assume $U$ and $V$ act irreducibly, and we denote the irreps as $\mu$ and $\nu$.

There is a convenient characterization of quantum channel covariance in terms of the Choi operator $\mathsf{C}_{\N}$, (cf. for instance, Lemma 11 in~\cite{gschwendtner2021programmability}.) 
\begin{equation}
    \text{$\N$ is $UV$-covariant}\ \Longleftrightarrow \ [\mathsf{C}_{\N},U^*\otimes V]=0\ .
\end{equation}
Here, the relevant representation is the dual (contragradient) representation $\mu^*$ of $\mu$ tensored with $\nu$. We can decompose $\mu^*\otimes\nu$ into irreps of U(d),
\begin{equation}\label{eq:decomp}
    \mu^*\otimes\nu \simeq \bigoplus_\lambda \lambda\ .
\end{equation}
Together with~\eqref{eq:covariance}, Schur's lemma implies that $\mathsf{C}_\N$ takes the form of a convex combination of projectors $\{\Pi_\lambda\}$ onto the irreps,
\begin{equation}\label{eq:covChoi}
    \mathsf{C}_\N\simeq \bigoplus_\lambda c_\lambda \frac{|\mu|}{|\lambda|} \Pi_\lambda\ , \quad \sum_\lambda c_\lambda =1, c_\lambda\ge 0\ ,
\end{equation}
where the factor $|\mu|/|\lambda|$ ensures that we have a convex combination of Choi operators that are normalized as $\tr\, (|\mu|/|\lambda|)C_\lambda=|\mu|$.

We are interested in a quantum channel that preserves the information as much as possible, so we naturally pick the Choi operator with the least (Kraus) rank. That is, we pick the projector onto the \emph{minimal irrep} in the above decomposition, denoted as $\omega$. Here minimality is defined with respect to the dominance/majorization (partial) order among the corresponding partitions of the irreps. While it is obvious that there is a maximal element $\mu^*+\nu$, there is also a unique minimum $\omega$ that is dominated/majorized by every other irrep $\lambda$ in the Littlewood-Richardson decomposition of $\mu^*\otimes\nu$, known as the \emph{Parthasarathy-Ranga Rao-Varadarajan (PRV)} component:
\begin{equation}
    \omega = (\mu^*+w_0(\nu))^+
\end{equation}
where $w_0$ is the longest Weyl group element that reflects the simple roots, $w_0:\alpha_i\mapsto -\alpha_{n-i}$, and the $(\cdot)^+$ denotes the projection onto the fundamental Weyl chamber so as to obtain a dominant weight. $\omega$ is then defined as the irrep with the highest weight $(\mu^*+w_0(\nu))^+$.

We can also write the dual as 
\begin{equation}
    \mu^* = - w_0(\mu)
\end{equation}
and,
\begin{equation}
    \omega =w_0(\nu-\mu)^+ = (\nu-\mu)^+
\end{equation}
that is,

In terms of the partitions or Young diagrams, $\omega$ is labeled by its highest weight, which is simply the difference $\nu-\mu$ re-ordered as a dominant weight.

We thence define our \emph{generalized cloning map} $\C_{\mu\to\nu}$ as
\begin{equation}\label{eq:gen_cloner}
\C_{\mu\to\nu}:\rho_\mu\mapsto\tr_\mu\,\mathsf{C}_{\C_{\mu\to\nu}}\cdot\rho^\top_\mu\otimes \mathbf{1}_\nu\ ,\quad \mathsf{C}_{\C_{\mu\to\nu}} :=\frac{|\mu|}{|\omega|} \Pi_{\omega}\ .
\end{equation}

The map $\C_{\mu\to\nu}$ is completely positive as its Choi operator is proportional to a projector. To show it is also trace-preserving, consider the linear function $f:\rho_\mu\mapsto\tr\,\C_{\mu\to\nu}(\rho_\mu)$. By covariance, we have 
\begin{equation}
    f(U_\mu(g)\rho_\mu U_\mu(g)^\dagger)=\tr\,(U_\nu(g)\C_{\mu\to\nu}(\rho_\mu)U_\nu(g)^\dagger)=\tr\,\C_{\mu\to\nu}(\rho_\mu)=f(\rho_\mu)\ .
\end{equation}
Since $\mu$ is an irrep, it admits a unique linear functional invariant under the group action, that is the trace (up to a prefactor). We have $f=\tr$, because of $\tr\,(\mathbf{1}_\mu/|\mu|) = \tr\, (\C_{\mu\to\nu}(\mathbf{1}_\mu/|\mu|))=\tr\,(\mathbf{1}_{\omega}/|\omega|)=1 $.

Let us do some quick sanity checks. When $\mu=\nu$, we have $\omega$ being the trivial irrep, and the projector on the trivial irrep corresponds to the identity channel. 

We can also examine the qubit case where the optimal cloner is known. All irreps are isomorphic to the symmetric subspaces Sym$^k(\mathbb{C}^2)$ for some $k$. For $\mu\simeq$ Sym$^n(\mathbb{C}^2)$ and $\nu\simeq$ Sym$^m(\mathbb{C}^2)$. The irreps-decomposition of $\mu^*\otimes\nu$ reads
\begin{equation}
    \mu^*\otimes\nu \simeq \bigoplus_{k=0}^{\min(m,n)} \mathrm{Sym}^{|n-m+2k|}(\mathbb{C}^2)\ .
\end{equation}
Hence, the minimal irrep is $\mathrm{Sym}^{|n-m|}(\mathbb{C}^2)$ which has the highest weight $(\nu-\mu)^+=(|n-m|,0)$. The Kraus rank of the corresponding generalized cloning channel is $|n-m|+1$, which uniquely matches the covariant channel used by YCH~\eqref{eq:qubitcloning}. Note that when $m<n$, the generalized cloning channel becomes a partial trace in this example.\\

Recall that Werner's map reads
\begin{equation}
    \mathcal{C}_{n \to m}(\rho) = \frac{d_{\mathrm{Sym}^n(\mathbb{C}^2)}}{d_{\mathrm{Sym}^m(\mathbb{C}^2)}}V^\dagger(\rho\otimes I_{n-m})V
\end{equation}
where $\rho$ is a state on Sym$^n(\mathbb{C}^2)$ and $I_{m-n}$ is the identity on $\mathrm{Sym}^{|n-m|}(\mathbb{C}^2)$, and $V$ is the isometric embedding of Sym$^m(\mathbb{C}^2)$ into Sym$^n(\mathbb{C}^2)\otimes \mathrm{Sym}^{|n-m|}(\mathbb{C}^2)$.\footnote{Usually the cloning map is written as $\mathcal{C}_{n \to m}(\rho) = \frac{d_{\mathrm{Sym}^n(\mathbb{C}^2)}}{d_{\mathrm{Sym}^m(\mathbb{C}^2)}}P_m(\rho\otimes I_{n-m})P_m$, where $P_m$ is the projection onto Sym$^m(\mathbb{C}^2)$ but $P_m$ is act on Sym$^n(\mathbb{C}^2)\otimes \mathrm{Sym}^{|n-m|}(\mathbb{C}^2)$. We take the intrinsic view by using the isometry $V$ instead, so the channel's codomain is strictly Sym$^m(\mathbb{C}^2)$.}

In fact, we can turn the Choi representation into its Stinespring representation that resembles Werner's cloning map more manifestly. Let $V: \mathcal{H}_\nu \to \mathcal{H}_\mu \otimes \mathcal{H}_\omega$ be the isometric embedding of the representation $\nu = \mu + \omega$ be the isometric embedding of the representation $\nu = \mu + \omega$ into $\mu\otimes\omega$. The generalized cloning map can be written as
\begin{equation}
    \mathcal{C}_{\mu \to \nu}(\rho) = \frac{d_\mu}{d_\nu} V^\dagger (\rho \otimes I_\omega) V\ .
\end{equation}
We justify our claim in the following lemma.
\begin{lem}[Choi Matrix of the Minimal Covariant Channel]\label{lem:choi_prv}
    Let $V: \mathcal{H}_\nu \to \mathcal{H}_\mu \otimes \mathcal{H}_\omega$ be the isometric embedding of the representation $\nu = \mu + \omega$. The covariant channel $\mathcal{C}_{\mu \to \nu}(\rho) = \frac{d_\mu}{d_\nu} V^\dagger (\rho \otimes I_\omega) V$ has a Choi matrix $\mathsf{C}$ that is proportional to the projector $\Pi_\omega$ onto the PRV component $(\nu - \mu)^+$ in $\nu \otimes \mu^*$. Specifically, $\mathsf{C} = \frac{d_\mu}{d_\omega} \Pi_\omega$.
\end{lem}

\begin{proof}
    The Choi matrix $\mathsf{C} \in \mathcal{B}(\mathcal{H}_\nu \otimes \mathcal{H}_{\mu^*})$ is defined by the action of the channel on the unnormalized maximally entangled state $\ket{\Omega} = \sum_{M} \ket{\mu, M} \otimes \ket{\mu^*, M^*}$:
    \begin{equation}
        \mathsf{C} = \frac{d_\mu}{d_\nu} \sum_{M, M'} \sum_W \left( V^\dagger \ketbra{\mu, M; \omega, W}{\mu, M'; \omega, W} V \right) \otimes \ketbra{\mu^*, M^*}{\mu^*, M'^*} \ .
    \end{equation}
    We define the map $\tilde{V}: \mathcal{H}_\nu \otimes \mathcal{H}_{\mu^*} \to \mathcal{H}_\omega$ by $\tilde{V} \ket{\nu, N; \mu^*, K^*} = \bra{\mu, K} V \ket{\nu, N}$. Evaluating the matrix elements of $\tilde{V}^\dagger \tilde{V}$, we find:
    \begin{equation}
        \bra{\nu, N; \mu^*, K^*} \tilde{V}^\dagger \tilde{V} \ket{\nu, N'; \mu^*, L^*} = \sum_W \bra{\nu, N} V^\dagger \ket{\mu, K; \omega, W} \bra{\mu, L; \omega, W} V \ket{\nu, N'} \ .
    \end{equation}
    Comparing this with the Choi matrix, we obtain the exact operator identity $\mathsf{C} = \frac{d_\mu}{d_\nu} \tilde{V}^\dagger \tilde{V}$.
    
    Because $V$ is an U$(d)$-covariant and so is $\tilde{V}$, by Schur's Lemma, $\tilde{V}^\dagger \tilde{V}$ must be proportional to the projection onto the $\omega$ isotypic component of $\nu \otimes \mu^*$. Since $\nu = \mu + \omega$, the representation $\omega$ is the unique highest weight PRV component $(\nu-\mu)^+$. Thus, $\tilde{V}^\dagger \tilde{V} \propto \Pi_\omega$.
    
    Finally, we evaluate the trace. Since $V$ is an isometry, $\Tr(\tilde{V}^\dagger \tilde{V}) = \sum_N \bra{\nu, N} V^\dagger V \ket{\nu, N} = d_\nu$. Because $\Tr(\Pi_\omega) = d_\omega$, we have $\tilde{V}^\dagger \tilde{V} = \frac{d_\nu}{d_\omega} \Pi_\omega$. The Choi matrix is therefore $\mathsf{C} = \frac{d_\mu}{d_\omega} \Pi_\omega$.
\end{proof}


Another variant of the generalized cloning map is to choose the smallest irrep according to the dimension of the irrep in the decomposition. Firstly, the smallest irrep is generally different from the minimal irrep we proposed to use.\footnote{One example is the U(4) irreps $(7, 5, 1, 1)$ and $(7, 4, 2, 1)$, which have dimensions $540$ and $630$ respectively. The former majorizes the latter, but the former has a smaller dimension.}. 

We propose to use the minimal irrep because, to our best knowledge, the smallest irrep does not admit a simple construction. It is plausible that the smallest irrep could also serve our purposes, but it is technically harder to work with. Furthermore, there are hints that the minimal irrep is a strictly better choice in some other universal information processing task where the covariant channel construction is used. For instance, in quantum purity amplification, the optimal protocol utilizes the generalized cloning map from any irrep $\mu$ to the fundamental irrep $\nu=\Box$. That is the Choi matrix of the map is proportional to the projection onto the minimal irrep $\omega\subset \mu^*\otimes\Box$, rather than the smallest irrep. Hence, this is good evidence for using the minimal irrep for the generalized cloning map between any two irreps. 

For a covariant endomorphism, the canonical way to assess the fidelity of the channel is to see how well the channel preserves any input state, either in the average-case scenario or worst-case scenario. For a covariant channel that maps between density operators on distinct irreps, we need a general way to assess the fidelity.

We now describe one such measure. The Schur-Weyl duality states that
\begin{equation}\label{eq:schur_d}
    \rho^{\otimes n}\simeq \bigoplus_\lambda p_\lambda \rho_\lambda \otimes \tau_\lambda,\quad (U(g)\rho U(g)^\dagger)^{\otimes n}\simeq \bigoplus_\lambda p_\lambda U_\lambda(g)\rho_\lambda U_\lambda(g)^\dagger \otimes \tau_\lambda\ .
\end{equation}
where $\rho_\lambda$ are the GL$(d)$ irreps of $\rho$. 

The eigenvectors of $\rho_\lambda$ in~\eqref{eq:schur_d} are labeled by \emph{weights} $w$ ($\sum_{i=1}^d w_i=n, w_i\in\mathbb{N}$). They are majorized by $\lambda$, denoted as $w\prec\lambda$, which diagrammatically means that the weights can be filled into $\lambda$ as a semi-standard Young tableau. Let the spectrum of $\rho$ be $r$. The eigenvalues are given by $r^w:=\prod_{i=1}^d r_i^{w_i}$ up to normalization. The weights are ordered lexicographically, which is equivalent to the ordering of the values $r^w$. The normalization factor is
\begin{equation}
    s_\lambda(r) =:\sum_{w \prec \lambda} m_{\lambda}(w)r^w = \frac{\sum_{\pi\in S_d} \mathrm{sgn}(\pi)\ \Pi_{i=1}^d p_i^{\lambda_{\pi(i)}+d-\pi(i)}}{\Pi_{i<j}(p_i-p_j)} \ ,
\end{equation}
where the degeneracy $m_{\lambda}(w):=K_{\lambda w}$ is the Kostka number; and $s_\lambda(r)$ is the character of the $\lambda$-irrep of GL(d,$\mathbb{C}$) that naturally extends U(d), also known as the \emph{Schur polynomials}. Hence, the eigenvalues of $\rho_\lambda$ are $\{r^w/\chi_\lambda(r)\}_w$. 


A natural way to gauge the channel fidelity is to see if it can preserve $\rho_\lambda$ as the GL$(d)$ irreps. We would like to show something like the following,
\begin{equation}\label{eq:cloning}
    \frac12||\C_{\mu\to \nu}(\rho_\mu) - \rho_\nu||_1\lesssim d(\mu,\nu)/n\ ,
\end{equation}
where $d(\mu,\nu):=\sum_{i=1}^d |\mu_i-\nu_i|$.

Specifically, since $p_\lambda$ concentrates around the most typical Young diagram $\lambda_{\mathbf{p}}=n\mathbf{p}$, and the width is of order $\O(\sqrt{n})$, we are only interested in
\begin{equation}\label{eq:error1}
    \frac12||\C_{\mu\to \lambda_{\mathbf{p}}}(\rho_\mu) - \rho_{\lambda_{\mathbf{p}}}||_1\lesssim \O(1/\sqrt{n}), \quad\forall\mu \ \mathrm{s.t.}\ d(\mu,\lambda_{\mathbf{p}})|\sim\sqrt{n}\ .
\end{equation}

For $d=2$, one has explicit formulae for all the U$(2)$ CG coefficients. This fact has been shown by YCH through direct calculations,
\begin{equation}
    \frac12||\C_{J\to K}(\rho_J) - \rho_K||_1\le |J-K|/J\sim\ |J-K|/n,\quad \forall J, K\ ,
\end{equation}
where subscript $J$ labels the Young diagram with partition $(n/2+J,n/2-J)$. 




\section{Estimating the generalizing cloning fidelity}
We prove a special case of~\eqref{eq:cloning} where $\mu$ dominates $\nu$ in the sense that $\omega=\nu-\mu$ is still an ordered partition, representing a valid Young diagram, but $\omega$ is much smaller than $\mu$ and $\nu$ in the $n$-scaling limit. This result could eventually help us establish that the asymptotically qudit-cloning fidelity is perfect if all we want is a sublinear amount of extra clones.

\begin{prop}\label{prop:main}
    Let $\rho$ be a density matrix on $\mathbb{C}^d$ with rank $r \le d$. We assume its non-zero eigenvalues are strictly non-degenerate, satisfying $x_1 > x_2 > \dots > x_r > 0$, while $x_{r+1} = \dots = x_d = 0$. Let $\mu,\nu,\omega$ be three ordered partitions labeling GL$(d)$ irreps such that $\mu+\omega=\nu$. Their sizes are parameterized by a large integer $n$ such that $|\mu|,|\nu|=\Theta(n)$ and $||\omega||=O(n^s)$ for some $0<s<1$. Let $\rho_{\mu}$ and $\rho_\nu$ be the normalized GL$(d)$ representations of $\rho$. Let $\C_{\mu\to\nu}$ be the generalized cloning map from $\mu$ to $\nu$, and similarly for $\C_{\nu\to\mu}$. For any arbitrarily small $\frac{1-s}{2}\epsilon>0$, the fidelity of the covariant recovery is asymptotically bounded by:
    \begin{equation}
        F(\C_{\mu\to\nu}(\rho_\mu), \rho_\nu)\ge 1-O(n^{s-1+\epsilon})\ ,\quad F(\C_{\nu\to\mu}(\rho_\nu), \rho_\mu)\ge 1-O(n^{s-1+\epsilon})\ .
    \end{equation}
\end{prop}

\subsection{Preliminaries}

In this section, we establish the representation-theoretic notation and algebraic conventions used throughout the paper. Let $G = \mathrm{GL}(d)$ and let $\mathfrak{g} = \mathfrak{gl}(d, \mathbb{C})$ be its complexified Lie algebra. We fix a standard Cartan subalgebra $\mathfrak{h}$ consisting of diagonal matrices, and denote the dual weight space by $\mathfrak{h}^*$. 

\paragraph{Roots, Coroots, and the Cartan-Weyl Basis}
The algebra decomposes into $\mathfrak{h}$ and the one-dimensional root spaces spanned by the root vectors $E_\alpha$. We choose the standard Cartan-Weyl basis $\{H_i, E_\alpha, E_{-\alpha}\}$, where $H_i \in \mathfrak{h}$ and $[E_\alpha, E_{-\alpha}] = H_\alpha$. 

We equip $\mathfrak{h}^*$ with the standard Euclidean inner product $(\cdot, \cdot)$ induced by the trace. In the canonical basis $\{e_i\}_{i=1}^d$ of $\mathbb{R}^d$, the positive roots are $\alpha = e_i - e_j$ for $i < j$. Let $\Phi^+$ denote the set of positive roots. The simple roots, which form a basis for the positive root lattice, are given by:
\begin{equation}
    \alpha_i = e_i - e_{i+1} \ , \quad i = 1, \dots, d-1 \ .
\end{equation}
and the root lattice $Q_+$ is defined as
\begin{equation}
    Q_+ = \left\{\sum_{i=1}^dc_i\alpha_i,\ \ c_i\in\mathbb Z_{\ge 0}\right\} .
\end{equation}


For any positive root $\alpha$, the triplet $\{E_\alpha, E_{-\alpha}, H_\alpha\}$ generates an $\mathfrak{sl}_2\{\alpha\}$ subalgebra. The abstract coroot $\alpha^\vee \in \mathfrak{h}$ is defined by its pairing with any weight $\lambda \in \mathfrak{h}^*$ as $\langle \lambda, \alpha^\vee \rangle = \frac{2(\lambda, \alpha)}{(\alpha, \alpha)}$. Because $\mathfrak{gl}(d)$ belongs to the simply-laced Type $A_{d-1}$ family, all roots have identical squared length $(\alpha, \alpha) = 2$. Consequently, the coroots are canonically identified with the roots themselves ($\alpha^\vee = \alpha$), and the pairing reduces strictly to the inner product: $\langle \lambda, \alpha^\vee \rangle = (\lambda, \alpha)$. 

For the simple roots, this pairing extracts the difference between adjacent Dynkin labels:
\begin{equation}
    \langle \lambda, \alpha_i^\vee \rangle = (\lambda, \alpha_i) = \lambda_i - \lambda_{i+1} \ .
\end{equation}

\paragraph{The Quadratic Casimir Operator}
The total quadratic Casimir operator $C_2$ commutes with the group action and acts as a scalar on any irreducible representation $V_\lambda$. In the Cartan-Weyl basis, it takes the form:
\begin{equation}
    C_2 = \sum_{i=1}^d H_i^2 + \sum_{\alpha\in\Phi^+} (E_{-\alpha} E_\alpha + E_\alpha E_{-\alpha}) \ .
\end{equation}
The exact Casimir eigenvalue on $V_\lambda$ is given by $C_2(\lambda) = (\lambda, \lambda + 2\rho)$, where $\rho = \frac{1}{2} \sum_{\alpha \in\Phi^+} \alpha$ is the Weyl vector. When analyzing tensor products $V_\mu \otimes V_\omega$, the total Casimir operator dictates the decomposition into irreducible subspaces, enabling tight spectral bounds via algebraic perturbation theory.

\paragraph{Weight Spaces and Root Paths}
Let $V_\mu$ denote the finite-dimensional irreducible representation of $\mathrm{GL}(d)$ uniquely labeled by its strictly dominant highest weight $\mu = (\mu_1, \dots, \mu_d)$, where $\mu_1 \ge \mu_2 \ge \dots \ge \mu_d$. 

The representation space decomposes into orthogonal weight spaces $W_\mu(\delta)$. A state in $W_\mu(\delta)$ has weight $\mu - \delta$, where the weight offset $\delta \in Q_+$ can be expanded in the simple root basis as $\delta = \sum_{k=1}^{d-1} c_k \alpha_k$ with non-negative integers $c_k$. 

We often describe this offset via a \emph{path} of simple roots descending from the highest weight. Let an ordered set $\vec\delta=(\alpha_{l_1},\ldots, \alpha_{l_T})$, where $l_i\in[d-1]$, denote such a path. We use $|\delta| = T$ to denote the length of the path (which matches the height $\sum c_k$). For an intermediate step $t$ along the path ($0 \le t \le |\delta|$), $\vec\delta_t$ denotes the ordered subset $(\alpha_{l_1},\ldots, \alpha_{l_t})$. 

The weight offset $w(\vec\delta_t)$ at step $t$ is exactly:
\begin{equation}
    w(\delta_t) = \sum_{i=1}^t\alpha_{l_i} \ , \quad 0 \leq t \leq T \ .
\end{equation}
Because multiple paths can lead to the exact same weight offset, we seek path-independent statements and will frequently use the unordered subset $\delta_t$ to denote the weight offset for brevity.

\paragraph{The Gelfand-Tsetlin Basis}
We work in the Gelfand-Tsetlin (GT) basis, which is highly convenient for explicit calculations. Each GT basis vector $\ket{M}$ is uniquely labeled by a GT pattern $M$, which is a triangular arrangement of non-negative integers:
\[
\begin{matrix}
x_{n1} && x_{n2} && \cdots && \cdots && x_{nn}\\
& \ddots && \ddots && && \iddots && \\
&& x_{21} && x_{22}\\
&&& x_{11}
\end{matrix}
\]

The entries, $M_{i,j}\ , \quad 1\leq i\leq d,\quad 1\leq j\leq i \ ,$ satisfy the interlacing constraint:
\begin{equation}
    M_{i,j} \ge M_{i-1,j} \ge M_{i,j+1} \ .
\end{equation}
The top row is fixed at $i=d$ and satisfies $M_{d,j} = \mu_j$. The lower rows are indexed in descending order. The $k$-th weight vector component of a pattern $M$ is defined as:
\begin{equation}
    w_k(M) \equiv \sum_{j=1}^k M_{k,j} - \sum_{j=1}^{k-1} M_{k-1,j} \ .
\end{equation}
We use $m_{i,j}$ to denote a standard basis element for GT patterns, representing a pattern containing $1$ at $(i,j)$ and zero elsewhere. 

\paragraph{Kostka Numbers and Weight Multiplicities}
The dimension of a given weight block $W_\mu(\delta)$ is given by the inner degeneracy, known as the Kostka number:
\begin{equation}
    m_\mu(\delta) = K_{\mu,\,\mu-\delta} \ , \quad m_\nu(\delta) = K_{\nu,\,\nu-\delta} \ .
\end{equation}
We rely on the following exact equivalence for these multiplicities in the shallow-depth regime:

\begin{lem}[Kostka numbers]\label{lem:kostka}
    Let $\mu,\omega$ be two $\mathrm{GL}(d)$ irreps, $\delta$ be a weight offset, and $g_\mu$ be the edge gap of $\mu$. The inner degeneracies satisfy:
    \begin{equation}
        K_{\mu,\mu-\delta} \leq K_{\mu+\omega,\mu+\omega-\delta} \ .
    \end{equation}
    Moreover, we have strict equality if $|\delta| \leq g_\mu$.
\end{lem}
\begin{proof}
We define the injection $\phi_\omega$ by specifying how it acts on the GT basis.
\begin{equation}
    \phi_\omega: M\mapsto M+\omega,\quad\text{where}\quad (M+\omega)_{i,j} \equiv M_{i,j} + \omega_j, \ \ 1\leq j\leq i\ .
\end{equation}
That is, $\phi_\omega$ adds $\omega$ to the GT pattern columnwise. 

To visualize this injection, consider the standard triangular arrangement of a GT pattern for $d=3$. The map $\phi_\omega$ adds the weight $\omega$ such that the same $\omega_j$ is added along the $j$-th diagonal (column):
\begin{equation}
    M = \begin{array}{ccccc}
        M_{3,1} & & M_{3,2} & & M_{3,3} \\
        & M_{2,1} & & M_{2,2} & \\
        & & M_{1,1} & & 
    \end{array}
    \quad \xmapsto{\ \phi_\omega\ } \quad
    M+\omega = \begin{array}{ccccc}
        M_{3,1}+\omega_1 & & M_{3,2}+\omega_2 & & M_{3,3}+\omega_3 \\
        & M_{2,1}+\omega_1 & & M_{2,2}+\omega_2 & \\
        & & M_{1,1}+\omega_1 & & 
    \end{array}
\end{equation}

Because the interlacing inequalities $M_{i,j} \ge M_{i-1,j} \ge M_{i,j+1}$ compare adjacent elements in the triangle, adding $\omega_j$ to the $j$-th diagonal trivially preserves the left inequality $(M_{i,j}+\omega_j \ge M_{i-1,j}+\omega_j)$. The right inequality also holds because $\omega_j\ge\omega_{j+1}$.

The top row is $M_{d,j}+\omega_j = (\mu+\omega)_j$, so $M+\omega$ is a valid GT pattern in the irrep $\mu+\omega$. Since $M$ has weights $\mu-\delta$, it remains to show that $M+\omega$ has weight $\mu+\omega-\delta$. We have 
\begin{equation}
    w_k(M+\omega)=\sum_{j=1}^k (M_{k,j}+ \omega_j) - \sum_{j=1}^{k-1} (M_{k-1,j} + \omega_j) =w_k(M)+ \omega_k =(\mu+\omega-\delta)_k\ .
\end{equation}
Therefore, the first claim follows.

For the second claim, we consider the inverse map $\psi_\omega: \ket{N}\mapsto \ket{N-\omega}$ defined by columnwise subtraction:
\begin{equation}
 (N-\omega)_{i,j} \equiv N_{i,j} - \omega_j \ .
\end{equation}
We must verify that for any valid GT pattern $N$ in the irrep $\mu+\omega$ with weight offset $\delta$ (where $|\delta|\leq  g_\mu$), the pattern $M = N-\omega$ is a valid GT pattern in $\mu$.

The top row condition is satisfied by definition: $M_{d,j} = (\mu+\omega)_j - \omega_j = \mu_j$. The weight condition $w(M) = \mu-\delta$ follows from the linearity of the weight function. The nontrivial constraint is the interlacing condition for the rows $i < d$:
\begin{equation}
 M_{i,j} \geq M_{i-1,j} \geq M_{i,j+1} \ .
\end{equation}
Substituting $M = N-\omega$, the left inequality $N_{i,j}-\omega_j \geq N_{i-1,j}-\omega_j$ holds trivially because $N$ is a valid pattern. The right inequality requires:
\begin{equation}
N_{i-1,j} - \omega_j \geq N_{i,j+1} - \omega_{j+1} \iff N_{i-1,j} - N_{i,j+1} \geq \omega_j - \omega_{j+1} \ .
\end{equation}

To show this holds, we relate the depth $|\delta|$ to the total deficit of GT entries from their highest weight values. 
Let $N_0$ be the highest GT pattern of shape $\nu=\mu+\omega$, so that $(N_0)_{k,j}=\nu_j$ for all $1\le j\le k\le d$. 
Define row sums $s_k(N)\equiv\sum_{j=1}^k N_{k,j}$. Recall that the weight components satisfy $w_k(N)=s_k(N)-s_{k-1}(N)$, hence
\begin{equation}
s_k(N)=\sum_{i=1}^k w_i(N).
\end{equation}
If $N\in W_\nu(\delta)$ has weight $w(N)=\nu-\delta$, then for each $k\le d-1$,
\begin{equation}
s_k(N)=\sum_{i=1}^k (\nu_i-\delta_i), 
\qquad 
s_k(N_0)=\sum_{i=1}^k \nu_i.
\end{equation}
Therefore,
\begin{equation}
\sum_{k=1}^{d-1}\sum_{j=1}^k\big(\nu_j-N_{k,j}\big)=\sum_{k=1}^{d-1}\big(s_k(N_0)-s_k(N)\big)=\sum_{k=1}^{d-1}\sum_{i=1}^k\delta_i =\sum_{i=1}^{d-1}(d-i)\delta_i.
\end{equation}
Writing $\delta=\sum_{l=1}^{d-1}n_l\alpha_l$, one has
$\delta_1=n_1$ and $\delta_i=n_i-n_{i-1}$ for $2\le i\le d-1$, so a telescoping sum gives
\begin{equation}
\sum_{i=1}^{d-1}(d-i)\delta_i=\sum_{l=1}^{d-1}n_l=|\delta|.
\end{equation}
In particular,
\begin{equation}
|\delta|=\sum_{k=1}^{d-1}\sum_{j=1}^k\big((\mu+\omega)_j-N_{k,j}\big),
\end{equation}
and since each summand is nonnegative, every individual deficit is bounded:
\begin{equation}
(\mu_j+\omega_j)-N_{k,j}\le |\delta| \quad \forall\,k,j,
\end{equation}
hence $N_{i-1,j}\ge (\mu_j+\omega_j)-|\delta|$.

Combining this with the trivial upper bound $N_{i,j+1} \leq (\mu+\omega)_{j+1}$, we can bound the gap in~$N$:
\begin{align}
N_{i-1,j} - N_{i,j+1} &\geq (\mu_j + \omega_j - |\delta|) - (\mu_{j+1} + \omega_{j+1}) = (\mu_j - \mu_{j+1}) + (\omega_j - \omega_{j+1}) - |\delta| \ .
\end{align}
The condition for $M$ to be valid ($N_{i-1,j} - N_{i,j+1} \geq \omega_j - \omega_{j+1}$) is thus satisfied if:
\begin{equation}
 (\mu_j - \mu_{j+1}) - |\delta| \geq 0 \ .
\end{equation}
Since $g_\mu = \min_j (\mu_j - \mu_{j+1})$, the condition $|\delta| \leq g_\mu$ ensures this inequality holds for all $j$. Thus, $\psi_\omega$ is a valid inverse to $\phi_\omega$, proving $K_{\mu,\mu-\delta} = K_{\mu+\omega,\mu+\omega-\delta}$.
\end{proof}

\paragraph{Representations and Asymptotic Parameters}
To control the error in our estimations, we define two key parameters for any partition or Young diagram $\lambda$. First, we define its \emph{width} as:
\begin{equation}
    ||\lambda|| \equiv \lambda_1 - \lambda_d \ .
\end{equation}
Since we conventionally set $\lambda_d=0$, $||\lambda||$ simply represents the largest element of the partition, corresponding to the width of its Young diagram. Second, we define its \emph{edge gap} as the minimal simple root pairing over the effective support. Because $\rho$ has rank $r \le d$, the representation theory effectively restricts to the subgroup $\mathrm{GL}(r)$. We define the edge gap over this support as:
\begin{equation}
    g_\lambda \equiv \min_{1\leq i\leq r-1}(\lambda_i-\lambda_{i+1}) \ .
\end{equation}
For typical representations in our scaling limit, the drop at the boundary $\lambda_r - 0 = \Theta(n)$ is also extensive, ensuring the gap is uniformly $\Theta(n)$ across the non-zero support.

For representations $\mu$ and $\omega$, our asymptotic bounds are controlled by the ratio:
\begin{equation}
    \xi \equiv \frac{||\omega||}{g_\mu} \ .
\end{equation}
We operate in the regime where $||\omega|| = o(g_\mu)$. Specifically, we consider $||\omega|| = O(n^s)$ for $0 < s < 1$ and strongly regular highest weights where $g_\mu \approx g_\nu = \Theta(n)$, ensuring that $\xi = o(1)$.

\subsection{Proof}

Let us now sketch the proof outline of Proposition~\ref{prop:main}. If we cast the channel $\C_{\mu\to\nu}$ in the Kraus representation, we observe that each Kraus operator corresponds to a vector $\ket{W}$ in the irrep $\omega$. This is obvious because the Choi operator projects to $\omega$. Our strategy is to simplify the problem by only estimating the contribution from the Kraus operator corresponding to the highest weight in $\omega$, and ignore the rest in evaluating the fidelity. We show in Lemma~\ref{lem:perturbation} that we do not lose much by dropping all the rest Kraus operators.

Furthermore, we truncate the lower weights in $\rho_\mu$ and $\rho_\nu$ and only consider the contributions from the high weights. The truncation error is estimated in Lemma~\ref{lem:tail}. Lastly, we need to estimate the classical contribution to the fidelity that is due to comparing the spectra of $\rho_\mu$ and $\rho_\nu$. This is analyzed in Lemma~\ref{lem:ratio}. With the additional help of Lemma~\ref{lem:kostka}, Lemma~\ref{lem:commutativity} and Lemma~\ref{lem:monotonicity}, we present the proof of Proposition~\ref{prop:main} in the end.

\begin{defn}[Asymptotical scaling regime]\label{def:scaling}
    We consider the scaling regime where the minimal simple root pairing satisfies $g_\mu \equiv \min_{1\le i \le r-1} \langle \mu, \alpha_i^\vee \rangle = \Theta(n)$, and the ancilla weight is bounded as $\|\omega\| = O(n^s)$ for some $0 < s < 1$. 
\end{defn}

\begin{lem}[Davis-Kahan Theorem~\cite{davis1970rotation}]\label{lem:davis_kahan}
    Let $H_0$ and $H = H_0 + V$ be Hermitian operators on a finite-dimensional Hilbert space $\h$. Let the spectrum of $H_0$ be partitioned into two disjoint sets $S_0$ and $S_1$, separated by a spectral gap $\Delta \equiv \min_{\lambda \in S_0, \lambda' \in S_1} |\lambda - \lambda'|$.
    Let $P_0$ be the orthogonal projector onto the eigenspace of $H_0$ corresponding to $S_0$, and let $P$ be the projector onto the corresponding perturbed eigenspace of $H$. 
    If the perturbation satisfies $\|V\| < \Delta$, the operator norm distance between the projectors is bounded by:
    \begin{equation}
        \|P - P_0\|=\sin\theta_{\max} \le \frac{\|V P_0\|}{\Delta - \|V\|} \ ,
    \end{equation}
    which is also the sine of the maximum principal angle $\theta_{\max}$ between the subspaces $\mathrm{im}(P)$ and $\mathrm{im}(P_0)$.
\end{lem}


\begin{lem}[Highest-Weight Subspace Perturbation]\label{lem:perturbation}
    Let $\mu$ and $\omega$ be irreps of GL$(d)$.  Let $P_\nu$ be the orthogonal projector onto the irrep $\nu = \mu + \omega$ within the tensor product $V_\mu \otimes V_\omega$. Let $P_0$ be the projector onto the unperturbed subspace $V_\mu \otimes \ketbra{W_0}{W_0}$, where $\ket{W_0}$ is the highest weight state of $\omega$. Let $\Pi_{\nu-\delta}$ be the projector onto the subspace $W_{\nu-\delta}$ of total weight $\nu - \delta$, where $\delta$ is the sum of simple roots with depth $|\delta| =\Theta(n^\epsilon)$ for some arbitrarily small $\epsilon>0$. Defining the restricted projectors $P_\nu^{(\delta)} \equiv \Pi_{\nu-\delta} P_\nu \Pi_{\nu-\delta}$ and $P_0^{(\delta)} \equiv \Pi_{\nu-\delta} P_0 \Pi_{\nu-\delta}$, their operator norm distance is bounded by:
    \begin{equation}
        \|P_\nu^{(\delta)} - P_0^{(\delta)}\| \le O\left(n^{\frac{s-1+\epsilon}{2}} \right) \ 
    \end{equation}
    in the scaling regime where $g_\mu = \Theta(n)$ and $\|\omega\| = O(n^s)$ for some $0 < s < 1$. 
\end{lem}

\begin{proof}
We work in type $A_{d-1}$ with simple roots $\alpha_k = e_k - e_{k+1}$ and coroots $\alpha_k^\vee = \alpha_k$. For a dominant highest weight $\lambda$, we use the standard pairing $\langle \lambda, \alpha_k^\vee \rangle = \lambda_k - \lambda_{k+1}$, and write $(\eta, \lambda) := \sum_k c_k \langle \lambda, \alpha_k^\vee \rangle$ when $\eta = \sum_k c_k \alpha_k$.

We restrict the global operators to the joint weight subspace $W_{\nu-\delta}$. The target representation $\nu$ is the eigenspace of the total quadratic Casimir operator $C_{\mathrm{tot}}$ corresponding to the eigenvalue $C_2(\nu)$. We decompose $C_{\mathrm{tot}} = H_0 + V$:
\begin{equation}
    H_0 = C_2(\mu)I_\mu \otimes I_\omega + I_\mu \otimes C_2(\omega) I_\omega + 2 \sum_{i=1}^d H_i^{(\mu)} \otimes H_i^{(\omega)} \ ,
\end{equation}
\begin{equation}
    V = 2 \sum_{\alpha > 0} \left( E_{-\alpha}^{(\mu)} \otimes E_\alpha^{(\omega)} + E_\alpha^{(\mu)} \otimes E_{-\alpha}^{(\omega)} \right) \ ,
\end{equation}
where $V$ is the perturbation to $H_0$. 

\paragraph{The Unperturbed Spectrum and Gap ($\Delta$) }
The subspace $W_{\nu-\delta}$ decomposes orthogonally into $\mathcal{H}_0^{(\delta)} \oplus \mathcal{H}_1^{(\delta)}$, where $\mathcal{H}_0^{(\delta)} = \mathrm{im}(P_0^{(\delta)})$.
For any state in $\mathcal{H}_0^{(\delta)}$, the operator $H_0$ acts diagonally with the exact eigenvalue:
\begin{equation}
    \lambda_0 = C_2(\mu) + C_2(\omega) + 2 (\mu - \delta, \omega) \ .
\end{equation}
For a state in $\mathcal{H}_1^{(\delta)}$, the $\omega$-register has weight $\omega - \eta$ for some positive root sum $\eta = \sum c_k \alpha_k > 0$. By weight conservation, the $\mu$-register has weight $\mu - \delta + \eta$. The eigenvalue is $\lambda_1(\eta)$, and the difference is exactly:
\begin{equation}
    \Delta E(\eta) = \lambda_0 - \lambda_1(\eta) = 2 (\eta, \mu) - 2 (\eta, \delta + \omega) + 2 (\eta, \eta) \ge 2 (\eta, \mu) - 2 \|\eta\| \, \|\delta + \omega\| \ .
\end{equation}
Using the simple root expansion, the dominant term is bounded uniformly:
\begin{equation}
    2(\eta, \mu) = 2 \sum_k c_k \langle \mu, \alpha_k^\vee \rangle \ge 2 g_\mu \sum_k c_k = 2 g_\mu \, |\eta| \ge 2 g_\mu \ .
\end{equation}
Using $\|\eta\| = O(|\delta|)$, $\|\delta + \omega\| \le O(|\delta|) + O(\|\omega\|)$, and our bounds $|\delta| =\Theta(n^\epsilon)$ and $\|\omega\| = O(n^s)$, we obtain a uniform lower bound on the gap:
\begin{equation}
    \Delta := \min_{\eta \neq 0} (\lambda_0 - \lambda_1(\eta)) \ge 2 g_\mu - O(|\delta| \, \|\omega\|) - O(|\delta|^2) = \Theta(n) \ .
\end{equation}

\paragraph{The Target Eigenvalue of $C_{\mathrm{tot}}$}
Using the exact $\mathfrak{sl}_d$ Casimir formula $C_2(\lambda) = (\lambda, \lambda + 2\rho)$, the highest possible representation in $\mu \otimes \omega$ is $\nu = \mu + \omega$ with eigenvalue $C_2(\nu)$.
Notice that $\lambda_0 - C_2(\nu) = -2(\delta, \omega) = O(|\delta| n^s) = o(n)$. Therefore, $C_2(\nu)$ lies well within the $o(n)$ neighborhood of $\lambda_0$.

If an irrep with highest weight $\nu - \gamma$ contributes to the restricted space $W_{\nu-\delta}$, then $\delta - \gamma$ is in the dominant root cone, so $|\gamma| \le |\delta|$.
For any such $\gamma \neq 0$, the Casimir eigenvalue is separated from $C_2(\nu)$ by:
\begin{equation}
    C_2(\nu) - C_2(\nu-\gamma) = (\nu, \nu + 2\rho) - (\nu-\gamma, \nu-\gamma + 2\rho) = 2(\gamma, \nu+\rho) - (\gamma, \gamma) \ .
\end{equation}
Evaluating this for $|\gamma| \le |\delta|$:
\begin{equation}
    2(\gamma, \nu+\rho) - (\gamma, \gamma) \ge 2(\gamma, \mu) - O(\|\gamma\| \, \|\omega\|) - O(\|\gamma\|) \ge 2g_\mu - o(n) = \Theta(n) \ .
\end{equation}
Consequently, $C_2(\nu)$ is the unique eigenvalue of $C_{\mathrm{tot}}$ in the $\Delta$-neighborhood of $\lambda_0$, confirming $P_\nu^{(\delta)}$ is the perturbed spectral projector required by Davis-Kahan.

\paragraph{Norm of the Effective Perturbation ($\|V P_0^{(\delta)}\|$)}
Because $|W_0\rangle$ is annihilated by all $E_\alpha^{(\omega)}$, the effective perturbation is strictly $V P_0^{(\delta)} = 2 \sum_{\alpha > 0} (E_\alpha^{(\mu)} \otimes E_{-\alpha}^{(\omega)}) P_0^{(\delta)}$.

The action of the ladder operators on the weight spaces is tightly constrained by the $\mathfrak{sl}_2\{\alpha\}$ subalgebras. A state at depth $\delta$ lies on an $\alpha$-weight string passing through the highest weight $\mu$. The total length of this string is exactly $m = \langle \mu, \alpha^\vee \rangle$, and the state's position along the string is bounded by $k = \langle \delta, \alpha^\vee \rangle$. Let $W_\mu(\le |\delta|) \equiv \bigoplus_{|\gamma| \le |\delta|} W_\mu(\gamma)$ denote the subspace of states up to depth $|\delta|$. The operator norm of the raising operator $E_\alpha^{(\mu)}$ restricted to this subspace is universally bounded by:
\begin{equation}
||E_\alpha^{(\mu)}||_{W_{\mu}(\le |\delta|)} \le \max_{0 \le k \le |\delta|} \sqrt{(k+1)(m-k)} \le \sqrt{(|\delta|+1)m} \ .
\end{equation}

Similarly, $\|E_{-\alpha}^{(\omega)} |W_0\rangle\| = \sqrt{\langle \omega, \alpha^\vee \rangle} \le O(n^{s/2})$. 
Summing over the finite number of positive roots yields:
\begin{equation}
    \|V P_0^{(\delta)}\| \le \sum_{\alpha > 0} 2 \|E_\alpha^{(\mu)}\| \cdot \|E_{-\alpha}^{(\omega)}\| \le O(\sqrt{|\delta| n}) \cdot O(n^{s/2}) = O\left( n^{\frac{1+s+\epsilon}{2}} \right) \ .
\end{equation}

\paragraph{Davis-Kahan Evaluation}
Applying the same $\mathfrak{sl}_2$ string arguments to the $\omega$-register at depth $\le |\delta|$ gives $\|V\|_{W_{\nu-\delta}} \le O(|\delta| n^{\frac{1+s}{2}})$. To ensure the perturbative regime $\|V\| < \Delta / 2$ is valid, we require:
\begin{equation}
    |\delta| \, n^{\frac{1+s}{2}} = o(n) \quad \Longleftrightarrow \quad |\delta| = o\left(n^{\frac{1-s}{2}}\right) \ .
\end{equation}
Our logarithmic cutoff $|\delta| =\Theta(n^\epsilon)$ comfortably satisfies this, yielding $\|V\| = o(\Delta)$.

Finally, because $|\delta| \ll g_\mu$, Lemma~\ref{lem:kostka} guarantees $K_{\mu, \mu-\delta} = K_{\nu, \nu-\delta}$. Thus, the restricted projectors satisfy the exact rank-matching condition $\mathrm{rank}(P_0^{(\delta)}) = \mathrm{rank}(P_\nu^{(\delta)})$. Applying the $\sin \Theta$ bound gives the operator distance:
\begin{equation}
    \|P_\nu^{(\delta)} - P_0^{(\delta)}\| \le \frac{2 \|V P_0^{(\delta)}\|}{\Delta} \le \frac{O\left( \sqrt{|\delta|} \, n^{\frac{1+s}{2}} \right)}{\Theta(n)} = O\left( n^{\frac{s-1+\epsilon}{2}} \right) \ .
\end{equation}
\end{proof}

The normalized GL$(d)$ irrep of $\rho$, $\rho_\mu\equiv\pi_\mu(\rho)/s_\mu(x)$, is block diagonalized in weights:
\begin{equation}
    \rho_\mu = \sum_{\delta\in Q_+} p_\mu(\delta) \Pi_{\mu-\delta}\ ,
\end{equation}
where we denote the weights using the difference $\delta$ from the highest weight $\mu$.

The eigenvalues of $\rho_\mu$ are
\begin{equation}
    p_\mu(\delta)\equiv\frac{x^{\mu-\delta}}{s_\mu(x)}\ ,
\end{equation}
where $s_{\mu}(x)\equiv\Tr\, \pi_{\mu}(\rho)$ is the Schur polynomial and $x=(x_i)_{i=1}^r$ is a shorthand for the non-zero spectral values of $\rho$; and we also use the shorthand $x^{w}\equiv\prod_{i=1}^r x_i^{w_i}$ for the weights $w=\mu-\delta$ and $\nu-\delta$.

We now control the error of truncating the lower weights. 
\begin{lem}[Tail mass]\label{lem:tail}
Let $\rho$ be a density matrix with rank $r \le d$ and strictly non-degenerate non-zero spectrum $x_1 > \dots > x_r > 0$. Let $\rho_\mu$ be its representation in the GL$(d)$ irrep $\mu$, with spectral decomposition $\rho_\mu = \sum_{\delta\in Q_+} p_\mu(\delta) \Pi_{\mu-\delta}$, where $\Pi_{\mu-\delta}$ projects onto the subspace of weight $\mu-\delta$.
The tail mass at depth $T$, defined as $\eps_\mu(T)\equiv\Tr\left(\bigoplus_{\delta:|\delta|>T} p_\mu(\delta)\Pi_{\mu-\delta}\right)=\sum_{\delta:|\delta|>T}p_\mu(\delta) m_\mu(\delta)$, is bounded by
\begin{equation}
\eps_\mu(T)\leq O\left(T^\kappa p^T\right)
\end{equation}
where $p \equiv \max_{1 \le i \le r-1} \{x_{i+1}/x_i \,|\, \mu_i > \mu_{i+1}\} \in (0,1)$ is the maximum relevant eigenvalue ratio, and $\kappa = r(r-1)/2 - 1$ is a constant depending on the effective dimension.
\end{lem}
\begin{proof}
The probability of a state with weight $\mu-\delta$ is bounded by the product of eigenvalue ratios corresponding to the simple roots in $\delta$. Specifically, lowering by a simple root $\alpha_i$ replaces one factor $x_i$ by $x_{i+1}$, multiplying the eigenvalues by $x_{i+1}/x_i\leq p$. Hence, at step $t$, 
\begin{equation}
    p_\mu(\delta_t)/p_\mu(\delta_0)=x^{\mu-\delta_t}/x^\mu\leq p^t  \ .  
\end{equation}
We have
\begin{equation}
    \eps_\mu(T) = \sum_{\delta:|\delta|>T}\frac{x^{\mu-\delta}}{s_\mu(x)}m_\mu(\delta)\leq\frac{x^\mu}{s_\mu(x)}\sum_{\delta:|\delta|>T} p^{|\delta|} m_\mu(\delta)\leq\sum_{t> T}\sum_{\delta:|\delta|=t} p^t m_\mu(\delta)\equiv \sum_{t> T} D(t)p^t
\end{equation}
where we used $x^\mu/s_\mu(x)\leq 1$ in the penultimate step, and $D(t)$ is the total degeneracy at depth $t$.

To bound $D(t)$, we note that the finite-dimensional irreducible representation $V_\mu$ is a quotient of the infinite-dimensional Verma module $M(\mu)$. Consequently, the inner multiplicity $m_\mu(\delta)$ is globally bounded above by the weight multiplicity of the Verma module. This Verma multiplicity is exactly Kostant's partition function, which counts the number of ways to write $\delta$ as a sum of positive roots. Since $\rho$ is supported only on the first $r$ dimensions, the allowed positive roots are exactly those within the effective $\mathrm{GL}(r)$ subspace:
\begin{equation}
    m_\mu(\delta) \le K(\delta) = \big| \big\{ n_\alpha \in \mathbb{Z}_{\ge 0} \ \big| \ \delta = \sum_{\alpha\in\Phi^+_{\mathrm{GL}(r)}}n_\alpha\, \alpha \big\} \big| \ .
\end{equation}
We can now take heights on both sides of the sum:
\begin{equation}
|\delta| = t = \sum_{\alpha \in \Phi^+_{\mathrm{GL}(r)}} n_\alpha |\alpha| \ .
\end{equation}
The total bound on the degeneracy $D(t)$ at a fixed depth $t$ corresponds to the number of non-negative integer solutions to this height equation. This is a partition problem of the integer $t$ into parts defined by the heights of the positive roots. 

The asymptotic scaling of the tail mass relies on interpreting the depth $t$ through partition theory. Schur's theorem states the following: Let $A = \{a_1, a_2, \dots, a_k\}$ be a set of positive integers such that $\gcd(a_1, \dots, a_k) = 1$. Let $p_A(n)$ denote the number of partitions of an integer $n$ into parts taken from the set $A$ (where parts may be repeated). Then, the asymptotic growth of the number of partitions as $n \to \infty$ is given by
\begin{equation}
p_A(n) \sim \frac{n^{k-1}}{(k-1)! \prod_{i=1}^k a_i} \ .
\end{equation}
The upper bound on the number of states at a weight depth $t$ is given by the number of ways to partition $t$ using the heights of the positive roots within the effective support, i.e., $A = \{|\alpha| \mid \alpha \in \Phi^+_{\mathrm{GL}(r)}\}$. Since the simple roots have unit height, the condition $\gcd(A)=1$ is satisfied. The number of parts is $k = |\Phi^+_{\mathrm{GL}(r)}| = r(r-1)/2$. Applying Schur's theorem, the total degeneracy bound $D(t)$ at depth $t$ scales as:
\begin{equation}
D(t) \sim t^{\frac{r(r-1)}{2} - 1} \ .
\end{equation}

Therefore, the maximum number of states at depth $t$ scales polynomially as $D(t) \sim t^\kappa$ with $\kappa = \frac{r(r-1)}{2} - 1$.

Substituting this scaling, the bound becomes:
\begin{equation}
    \eps_\mu(T) \leq C \sum_{t > T} t^\kappa p^t \ .
\end{equation}
Since $p < 1$, the sum is dominated by the term at $t=T$ (multiplied by a geometric series factor), yielding $\eps_\mu(T) \leq O(T^\kappa p^T)$.
\end{proof}

Finally, we control the discrepancy between the spectra of $\rho_\mu$ and $\rho_\nu$. They can be spectrally decomposed as
\begin{equation}
    \rho_\mu = \sum_{\delta\in Q_+} p_\mu(\delta) \Pi_{\mu-\delta},\quad \rho_\nu = \sum_{\delta\in Q_+} p_\nu(\delta) \Pi_{\nu-\delta}
\end{equation}
Because of $\mu+\omega=\nu$, we have a $\delta$-independent ratio of $p_\nu(\delta)/p_\mu(\delta)$,
\begin{equation}
    Z(x)\equiv\frac{p_\nu(\delta)}{p_\mu(\delta)} = \frac{s_\mu(x)}{s_\nu(x)}x^\omega\ .
\end{equation}
Now we show that $Z(x)\approx 1$.

\begin{lem}[Probability ratio]\label{lem:ratio}
    Let $\mu,\omega$ be two GL$(d)$ irreps, and let $x$ be a probability vector with $r \le d$ strictly decreasing non-zero entries $x_1>\cdots>x_r>0$. Let $\delta_x \equiv \min_{1 \le i \le r-1} (\ln x_i - \ln x_{i+1})$ be the logarithmic spectral gap. The ratio $Z(x)\equiv x^\omega s_\mu(x)/s_{\mu+\omega}(x)$ satisfies
    \begin{equation}
        1 \ge Z(x) \ge 1 - O\left(e^{-(  g_\mu+1)\delta_x}\right)\ .
    \end{equation}
\end{lem}

\begin{proof}
    We first show the upper bound $Z(x)\leq 1$. This holds generally for any dominant integral weights. Using the expansion of Schur polynomials into standard Young tableaux (SSYT),
    \begin{equation}
        s_\mu(x) = \sum_{u\prec\mu} K_{\mu,u}\,x^u\ .
    \end{equation}
    Multiplying by $x^\omega$ shifts the weights: $x^\omega s_\mu(x) = \sum_{u} K_{\mu,u}\,x^{u+\omega}$. By Lemma~\ref{lem:kostka}, the Kostka numbers satisfy $K_{\mu,u} \leq K_{\mu+\omega, u+\omega}$. Since the weights of $\mu+\omega$ contain the shifted weights of $\mu$ as a subset, every term in the expansion of $x^\omega s_\mu(x)$ appears in $s_{\mu+\omega}(x)$ with a coefficient at least as large. Thus $s_{\mu+\omega}(x) - x^\omega s_\mu(x) \ge 0$, implying $Z(x) \leq 1$.

    For the lower bound, because $x$ has exactly $r$ non-zero entries, the Schur polynomials evaluate identically to those of $\mathrm{GL}(r)$. We utilize the determinantal formula for the Schur polynomial over these $r$ variables, $s_\lambda(x) = \det(x_i^{\lambda_j+r-j}) / \det(x_i^{r-j})$. The denominators cancel in the ratio $Z(x)$, leaving
    \begin{equation}
        Z(x) = x^\omega \frac{\det(x_i^{\ell_j})}{\det(x_i^{L_j})} \ ,
    \end{equation}
    where $\ell_j = \mu_j + r - j$ and $L_j = (\mu+\omega)_j + r - j = \ell_j + \omega_j$. We expand the determinants as sums over the symmetric group $S_r$:
    \begin{equation}
        \det(x_i^{\ell_j}) = \sum_{\sigma\in S_r} \mathrm{sgn}(\sigma) \prod_{j=1}^r x_j^{\ell_{\sigma(j)}} = \sum_{\sigma\in S_r} \mathrm{sgn}(\sigma) e^{\sum_{j} \ell_{\sigma(j)} \ln x_j} \ .
    \end{equation}
    Since both $\ell$ and $\ln x$ are strictly decreasing sequences, the sum is dominated by the identity permutation $\sigma = \mathrm{id}$. The exponent for the identity is $E_{\mathrm{id}} = \sum \ell_j \ln x_j$. For any other permutation $\sigma$, the exponent is strictly smaller.
    
    Consider a permutation $\sigma$ and its decomposition into adjacent transpositions. For a single adjacent swap $\tau=(k, k+1)$, the change in the exponent is
    \begin{equation}
        \Delta E = (\ell_{k+1}\ln x_k + \ell_k \ln x_{k+1}) - (\ell_k \ln x_k + \ell_{k+1} \ln x_{k+1}) = -(\ell_k - \ell_{k+1})(\ln x_k - \ln x_{k+1}) \ .
    \end{equation}
    Note that $\ell_k - \ell_{k+1} = \mu_k - \mu_{k+1} + 1 \ge   g_\mu + 1$. Thus, the gap is bounded by $\epsilon_\mu \equiv (  g_\mu+1)\delta_x$.
    The contribution of non-identity terms is suppressed exponentially:
    \begin{equation}
        \det(x_i^{\ell_j}) = e^{E_{\mathrm{id}}} \left( 1 + \sum_{\sigma\neq \mathrm{id}} \mathrm{sgn}(\sigma) e^{-(E_{\mathrm{id}} - E_\sigma)} \right) \ge x^\ell \left( 1 - (r!-1)e^{-\epsilon_\mu} \right) \ .
    \end{equation}
    Similarly, for the denominator with weights $L$, the gap is controlled by $\mu+\omega$:
    \begin{equation}
        \det(x_i^{L_j}) \le x^L \left( 1 + (r!-1)e^{-\epsilon_{\mu+\omega}} \right) \ .
    \end{equation}
    Since $\omega$ is a dominant weight, $\Delta_{\mu+\omega} \ge   g_\mu$, and thus $\epsilon_{\mu+\omega} \ge \epsilon_\mu$. Combining these estimates:
    \begin{equation}
        Z(x) \ge x^\omega \frac{x^\ell (1 - (r!-1)e^{-\epsilon_\mu})}{x^L (1 + (r!-1)e^{-\epsilon_\mu})} = \frac{1 - (r!-1)e^{-\epsilon_\mu}}{1 + (r!-1)e^{-\epsilon_\mu}} \ .
    \end{equation}
    Using $(1-y)/(1+y) \ge 1-2y$, we obtain $Z(x) \ge 1 - O(e^{-(  g_\mu+1)\delta_x})$.
\end{proof}

Finally, we need two lemmas that help us simplify the fidelity $F(\C_{\mu\to\nu}(\rho_\mu),\rho_\nu)$.

\begin{lem}[Commutativity]\label{lem:commutativity}
Let $\rho \in \mathcal{D}(\mathbb{C}^d)$ be a density matrix. Let $\mu$ and $\nu$ be irreps of $\mathrm{GL}(d)$, $ \rho_\mu$ and $ \rho_\nu$ be the normalized GL$(d)$ irreps of $\rho$.
Let $\C_{\mu \to \nu}: \mathcal{L}(V_\mu) \to \mathcal{L}(V_\nu)$ be a quantum channel that is U$(d)$-covariant, i.e., for all $g \in \text{U}(d)$:
\begin{equation}
    \C_{\mu \to \nu}( U_\mu(g) X U_\mu(g)^{-1} ) = U_\nu(g) \C_{\mu \to \nu}(X) U_\nu(g)^{-1}
\end{equation}
where $U_\mu(g) \equiv \pi_\mu(g)$ for $g \in \text{U}(d)$.
Then, we have
\begin{equation}
    \left[ \C_{\mu \to \nu}(\rho_\mu), \rho_\nu \right] = 0 \ .
\end{equation}
\end{lem}

\begin{proof}
The commutation relation is basis-independent. We work in the eigenbasis of the seed state $\rho$. Let $V = \mathbb{I}$, such that $\rho = \Lambda$ is a diagonal matrix. Consequently, the representation states $\rho_\mu$ and $\rho_\nu$ are images of the diagonal matrix $\Lambda$:
\begin{equation}
    \rho_\mu \propto \pi_\mu(\Lambda), \quad \rho_\nu \propto \pi_\nu(\Lambda)
\end{equation}
These operators are diagonal in the GT basis of their respective representation spaces. 

Let $T \subset \text{U}(d)$ be the maximal torus consisting of diagonal unitary matrices with determinant 1. Since $\Lambda$ is diagonal, it commutes with all diagonal matrices, including those in~$T$:
\begin{equation}
    [\rho_\mu, U_\mu(t)] = 0 \quad \forall t \in T
\end{equation}

Using the $\text{U}(d)$-covariance property for an element $t \in T$:
\begin{equation}
    \C_{\mu\to\nu}( U_\mu(t) \rho_\mu U_\mu(t)^{-1} ) = U_\nu(t) \C_{\mu\to\nu}(\rho_\mu) U_\nu(t)^{-1}
\end{equation}
Since $\rho_\mu$ commutes with $U_\mu(t)$, the left-hand side simplifies to $\C_{\mu\to\nu}(\rho_\mu)$. Rearranging yields:
\begin{equation}
    [\C_{\mu\to\nu}(\rho_\mu), U_\nu(t)] = 0 \quad \forall t \in T
\end{equation}

The representation space $V_\nu$ decomposes into weight spaces $V_\nu = \bigoplus_{w} V_\nu(w)$.  The weights $w$ in a fixed irrep of $\text{U}(d)$ are distinct even when restricted to the traceless Cartan subalgebra $\mathfrak{sl}_d$, which generated $T$. Thus, the $\text{U}(d)$ torus action distinguishes these subspaces via unique characters:
\begin{equation}
    U_\nu(t) \ket{v_w} = w(t) \ket{v_w}
\end{equation}
Since $\C_{\mu\to\nu}(\rho_\mu)$ commutes with the entire torus $T$, by Schur's Lemma, it must be block-diagonal with respect to the weight spaces:
\begin{equation}
    \C_{\mu\to\nu}(\rho_\mu) = \bigoplus_{w} B_w, \quad \text{where } B_w \in \mathcal{L}(V_\nu(w))
\end{equation}

The target state $\rho_\nu \propto \pi_\nu(\Lambda)$ is the image of a diagonal matrix. It acts as a scalar multiple of the identity on each weight space:
\begin{equation}
    \rho_\nu |_{V_\nu(w)} = (\textstyle\prod_i \lambda_i^{w_i}) \mathbb{I}_{V_\nu(w)}
\end{equation}
Since $\C_{\mu\to\nu}(\rho_\mu)$ is block-diagonal and $\rho_\nu$ is block-scalar, they commute:
\begin{equation}
    [\C_{\mu\to\nu}(\rho_\mu), \rho_\nu] = 0
\end{equation}
\end{proof}

\begin{lem}[Monotonicity]\label{lem:monotonicity}
Let $\mathcal{N}$ be a quantum channel and $\rho = \mathcal{N}(\eta)$ for some input state $\eta$. Let $K$ be a Kraus operator of $\mathcal{N}$ such that $\mathcal{N}(\cdot) = K(\cdot)K^\dagger + \sum_{i} L_i (\cdot) L_i^\dagger$, and define the unnormalized state $\rho' = K \eta K^\dagger$. If $\sigma$ is a density operator that commutes with $\rho$, i.e., $[\rho, \sigma] = 0$, then the fidelity satisfies:
\begin{equation}
    F(\rho', \sigma)\leq F(\rho, \sigma) \ .
\end{equation}
\end{lem}
Note that $\rho'$ is sub-normalized, but we still use the standard definition for fidelity in $F(\rho', \sigma)$ rather than the generlized fidelity for sub-normalized states. 

\begin{proof}
Let the spectral decomposition of $\sigma$ be $\sigma = \sum_k s_k |k\rangle\langle k|$. Since $[\rho, \sigma] = 0$, $\rho$ is diagonal in the same basis, given by $\rho = \sum_k x_k |k\rangle\langle k|$, where $x_k = \langle k | \rho | k \rangle$.

Let $\mathcal{D}_\sigma$ denote the dephasing channel in the eigenbasis of $\sigma$, defined as 
\begin{equation}
    \mathcal{D}_\sigma(X) = \sum_k \langle k | X | k \rangle |k\rangle\langle k|\ .
\end{equation}
By the monotonicity of fidelity under quantum channels, we have:
\begin{equation}
    F(\rho', \sigma) \leq F(\mathcal{D}_\sigma(\rho'), \mathcal{D}_\sigma(\sigma)).
\end{equation}
Since $\sigma$ is diagonal in this basis, $\mathcal{D}_\sigma(\sigma) = \sigma$, so
\begin{equation}
    F(\mathcal{D}_\sigma(\rho'), \mathcal{D}_\sigma(\sigma)) = F(\mathcal{D}_\sigma(\rho'), \sigma)\ .
\end{equation}
The state $\mathcal{D}_\sigma(\rho')$ is diagonal with entries $x'_k = \langle k | \rho' | k \rangle$. For commuting states, the fidelity reduces to the classical fidelity of their eigenvalue distributions:
\begin{equation}
    F(\mathcal{D}_\sigma(\rho'), \sigma) = \sum_k \sqrt{x'_k s_k}.
\end{equation}
By the definition of the channel $\mathcal{N}$, we can decompose $\rho$ as $\rho = \rho' + \rho_{\text{rem}}$, where $\rho_{\text{rem}} = \sum_i L_i \eta L_i^\dagger$. Since $\eta \geq 0$, $\rho_{\text{rem}}$ is a positive operator. Therefore, the diagonal elements satisfy:
\begin{equation}
    x_k = \langle k | \rho' | k \rangle + \langle k | \rho_{\text{rem}} | k \rangle = x'_k + \delta_k,
\end{equation}
where $\delta_k = \langle k | \rho_{\text{rem}} | k \rangle \geq 0$. This implies $x_k \geq x'_k$ for all $k$. Since the square root function is monotonic and $s_k \geq 0$, we have:
\begin{equation}
    \sum_k \sqrt{x'_k s_k} \leq \sum_k \sqrt{x_k s_k}.
\end{equation}
which is
\begin{equation}
    F(\mathcal{D}_\sigma(\rho'), \sigma) \leq  F(\rho, \sigma).
\end{equation}
Combining these inequalities, we conclude $F(\rho', \sigma) \leq F(\rho, \sigma)$.
\end{proof}

Now we are ready to prove the main proposition.

\begin{proof}[Proof of Proposition~\ref{prop:main}]
We first prove the bound for the forward channel $\mathcal{C}_{\mu\to\nu}$. To facilitate geometric comparison using the Davis-Kahan bounds from Lemma~\ref{lem:perturbation}, we identify the target space $\mathcal{H}_\nu$ with its isometric embedding inside the tensor product $\mathcal{H}_\mu \otimes \mathcal{H}_\omega$. Consequently, we write the generalized cloning map using the orthogonal projector $P_\nu$ rather than the isometry~$V$.

\begin{equation}
    \C_{\mu\to\nu}(\rho_\mu) = \frac{d_\mu}{d_\nu} P_\nu \left( \rho_\mu \otimes I_\omega \right) P_\nu \ .
\end{equation}
Because $VV^\dagger = P_\nu$ and fidelity is invariant under isometry, this modification does not change the fidelity we want to estimate.

We consider the specific Kraus operator corresponding to the highest-weight pure state in the ancilla space. This completely positive trace-non-increasing map, denoted $\mathcal{K}_\omega$, is geometrically equivalent to appending the ancilla in its highest-weight state $|W_0\rangle$ and projecting the joint system onto the irreducible $\nu$ subspace using the exact target projector $P_\nu$:
\begin{equation}\label{eq:kraus_proj}
    \mathcal{K}_\omega(\rho_\mu) = \frac{d_\mu}{d_\nu} P_\nu \left( \rho_\mu \otimes \ketbra{W_0}{W_0}_\omega \right) P_\nu \ .
\end{equation}
The states $\mathcal{C}_{\mu\to\nu}(\rho_\mu)$ and $\rho_\nu$ commute by Lemma~\ref{lem:commutativity}. Using the monotonicity of fidelity as in Lemma~\ref{lem:monotonicity}, we lower bound the fidelity of the full channel by that of the restricted map $\mathcal{K}_\omega$:
\begin{equation}
    F(\mathcal{C}_{\mu\to\nu}(\rho_\mu), \rho_\nu) \ge F(\mathcal{K}_\omega(\rho_\mu), \rho_\nu) = F\left(\bigoplus_\delta p_\mu(\delta)\mathcal{K}_\omega(\Pi_{\mu-\delta}), \bigoplus_{\delta} p_\nu(\delta)\Pi_{\nu-\delta}\right) \ .
\end{equation}
Since $\mathcal{K}_\omega$ maps the weight space $W_\mu(\delta)$ strictly into $W_\nu(\delta)$, we evaluate the fidelity block-by-block:
\begin{equation}\label{eq:fidelity1}
   F(\mathcal{C}_{\mu\to\nu}(\rho_\mu), \rho_\nu) \ge \sum_{\delta\in Q_+} \sqrt{p_\mu(\delta)p_\nu(\delta)} \cdot F(\mathcal{K}_\omega(\Pi_{\mu-\delta}), \Pi_{\nu-\delta}) \ .
\end{equation}

\paragraph{Fidelity of weight blocks via Subspace Geometry}
We evaluate the fidelity for a fixed weight block $\delta$:
\begin{equation}
    F_{\mathrm{block}}(\delta) \equiv F(\mathcal{K}_\omega(\Pi_{\mu-\delta}), \Pi_{\nu-\delta}) = \mathrm{Tr} \sqrt{ \Pi_{\nu-\delta} \mathcal{K}_\omega(\Pi_{\mu-\delta}) \Pi_{\nu-\delta} } \ .
\end{equation}
Substituting~\eqref{eq:kraus_proj}, the operator within the square root becomes $\frac{d_\mu}{d_\nu} P_\nu (\Pi_{\mu-\delta} \otimes \ketbra{W_0}{W_0}) P_\nu$.
We identify the operator $P_0^{(\delta)} \equiv \Pi_{\mu-\delta} \otimes \ketbra{W_0}{W_0}$ as the unperturbed projector restricted to the total weight $\nu-\delta$ from Lemma~\ref{lem:perturbation}. Similarly, let $P_\nu^{(\delta)} \equiv \Pi_{\nu-\delta} P_\nu \Pi_{\nu-\delta}$ be the exact target projector restricted to the same weight. 
The block fidelity simplifies to the trace of a product of restricted projectors:
\begin{equation}
    F_{\mathrm{block}}(\delta) = \sqrt{\frac{d_\mu}{d_\nu}} \mathrm{Tr} \sqrt{ P_\nu^{(\delta)} P_0^{(\delta)} P_\nu^{(\delta)} } \ .
\end{equation}
By the spectral properties of subspace projections, the non-zero eigenvalues of $P_\nu^{(\delta)} P_0^{(\delta)} P_\nu^{(\delta)}$ are $\cos^2 \theta_i$, where $\theta_i$ are the principal angles between the subspaces $\mathrm{im}(P_\nu^{(\delta)})$ and $\mathrm{im}(P_0^{(\delta)})$. The trace of the square root is the sum of their cosines:
\begin{equation}
    F_{\mathrm{block}}(\delta) = \sqrt{\frac{d_\mu}{d_\nu}} \sum_{i=1}^{m(\delta)} \cos \theta_i \ ,
\end{equation}
where the number of non-zero principal angles is exactly the subspace rank, which by Lemma~\ref{lem:kostka} is the weight multiplicity $m(\delta)$.

\paragraph{Perturbation estimates}
We bound the sum of cosines using the maximum principal angle~$\theta_{\max}$. Using the elementary inequality $\cos \theta \ge 1 - \sin^2 \theta$, we have $\cos \theta_i \ge 1 - \sin^2 \theta_{\max}$. 
The sine of the maximum principal angle is precisely the operator distance between the two projectors. Applying the Davis-Kahan bound from Lemma~\ref{lem:perturbation}:
\begin{equation}
    \sin \theta_{\max} = \|P_\nu^{(\delta)} - P_0^{(\delta)}\| \le O\left(\sqrt{|\delta|} \, n^{\frac{s-1}{2}}\right) \ .
\end{equation}
Squaring this bound, we obtain:
\begin{equation}
    \sum_{i=1}^{m(\delta)} \cos \theta_i \ge m(\delta) \left( 1 - \|P_\nu^{(\delta)} - P_0^{(\delta)}\|^2 \right) \ge m(\delta) \left( 1 - O(|\delta| n^{s-1}) \right) \ .
\end{equation}
Substituting this back, the block fidelity is rigorously bounded by:
\begin{equation}
    F_{\mathrm{block}}(\delta) \ge \sqrt{\frac{d_\mu}{d_\nu}} \, m(\delta) \left( 1 - O(|\delta| n^{s-1}) \right) \ .
\end{equation}
We set a weight offset cutoff at $|\delta|\le T = O(n^\epsilon)$ for an arbitrarily small $\epsilon > 0$. For states at depth less than that, the subspace deviation error is bounded uniformly by $O(n^{s-1+\epsilon})$. 

\paragraph{Dimension ratio}
Using the Weyl Dimension Formula, the ratio of dimensions over the full space is:
\begin{equation}
    \frac{d_\mu}{d_\nu} = \prod_{1 \le i < j \le d} \frac{\mu_i - \mu_j + j - i}{\nu_i - \nu_j + j - i} \ .
\end{equation}
Substituting $\nu_i = \mu_i + \eta_i$ and assuming $n$ is large such that the shift $j-i$ is negligible compared to row differences:
\begin{equation}
    \frac{d_\mu}{d_\nu} \approx \prod_{1 \le i < j \le d} \left( 1 + \frac{\eta_i - \eta_j}{\mu_i - \mu_j} \right)^{-1} \ .
\end{equation}
Taking the logarithm and using the Taylor expansion $\ln(1+x) \approx x$ for small $x \sim O(n^{s-1})$:
\begin{equation}
    \ln \left( \frac{d_\mu}{d_\nu} \right) \approx - \sum_{1 \le i < j \le d} \frac{\eta_i - \eta_j}{\mu_i - \mu_j} \ .
\end{equation}
We can rewrite the sum over pairs as a sum over rows acting against the ``potential'' generated by other rows. Because $\mu_k = \nu_k = 0$ for $k > r$, the numerator $\eta_k = 0$ for $k > r$. Thus, the non-zero contributions strictly arise when $k \le r$:
\begin{equation}
    \sum_{i < j} \frac{\eta_i - \eta_j}{\mu_i - \mu_j} = \sum_{k=1}^d \eta_k \sum_{j \neq k} \frac{1}{\mu_k - \mu_j} = \sum_{k=1}^r \eta_k \sum_{j \neq k} \frac{1}{\mu_k - \mu_j} \ .
\end{equation}
Because the edge gap across the support scales as $g_\mu = \Theta(n)$ and the boundary drop $\mu_r - 0 = \Theta(n)$, every term in the denominator $\mu_k - \mu_j$ for $k \le r$ scales as $\Theta(n)$. This ensures that the sum over the inverse differences is strictly bounded by order $O(n^{-1})$. Exponentiating back, and noting that $\nu_k - \mu_k = \eta_k = O(n^s)$, we obtain the leading order asymptotic behavior:
\begin{equation}
    \frac{d_\mu}{d_\nu} = 1 - \sum_{k=1}^r (\nu_k - \mu_k) \sum_{j \neq k} \frac{1}{\mu_k - \mu_j} + O(n^{2(s-1)})= 1-O(n^{s-1}) \ .
\end{equation}

\paragraph{Global fidelity and depth cutoff}
Substituting the perturbation estimate and the dimension ratio into Eq.~\eqref{eq:fidelity1}:
\begin{align}
    F(\mathcal{C}_{\mu\to\nu}(\rho_\mu), \rho_\nu) &\ge \sum_{|\delta|\le T} \sqrt{p_\mu(\delta) p_\nu(\delta)} \, m(\delta) (1 - O(n^{s-1+\epsilon})) \nonumber \\
    &\ge \left(1 - O(n^{s-1+\epsilon})\right) \sum_{|\delta|\le T} p_\mu(\delta) m(\delta) \sqrt{\frac{p_\nu(\delta)}{p_\mu(\delta)}} \ .
\end{align}
From Lemma~\ref{lem:ratio}, the probability ratio is uniformly $1 - O(e^{-n})$.
Finally, we account for the truncated tail mass. By Lemma~\ref{lem:tail}, the omitted probability mass beyond depth $T$ decays exponentially as $O(T^{\kappa} p^T)$. With our choice of $T = O(n^\epsilon)$, this error scales as $O(n^{\kappa\epsilon}p^{n^\epsilon})$ which is superpolynomially small. Hence, we can safely remove the constraint in the sum and obtain $\sum_{\delta\in Q_+} p_\mu(\delta) m(\delta)=1$ by definition.

Combining all bounds, the fidelity is constrained strictly by the subspace deviation error:
\begin{equation}
    F(\mathcal{C}_{\mu\to\nu}(\rho_\mu), \rho_\nu) \ge 1 - O(n^{s-1+\epsilon}) \ .
\end{equation}

\paragraph{Reverse Channel $\mathcal{C}_{\nu\to\mu}$}
The generalized cloning map from $\nu$ to $\mu$ uses the conjugate representation $\omega^* = (\mu-\nu)^+$. The geometric derivation parallels the forward case exactly, identifying $\mu$ as the target representation inside $\nu \otimes \omega^*$. Because the map goes from $\nu$ to $\mu$, the trace-preserving prefactor is inverted, evaluating exactly to $\sqrt{d_\nu / d_\mu} = (1 - O(n^{s-1}))^{-1/2} = 1 + O(n^{s-1})$.

Multiplying this prefactor by the Davis-Kahan subspace deviation gives a block fidelity bounded by $(1 + O(n^{s-1}))(1 - O(|\delta| n^{s-1})) \ge 1 - O(|\delta| n^{s-1})$. The dominant error remains the geometric subspace deviation $O(|\delta| n^{s-1})$, ensuring the identical asymptotic bound holds:
\begin{equation}
    F(\mathcal{C}_{\nu\to\mu}(\rho_\nu), \rho_\mu) \ge 1 - O(n^{s-1+\epsilon}) \ .
\end{equation}
\end{proof}


\section{Achievability for qudits}
The achievability upper bound reads
\begin{equation}\label{eq:achievable_degenerate}
    \log|M|\le \frac{d(d-1)-\sum_i g_i(g_i-1)}{2}\log n+\sum_{i<j}g_ig_j\log(p_i-p_j) -\sum_{k=1}^{d-1}\log k!+\sum_i\sum_{k=1}^{g_i}\log k!\ ,
\end{equation}



\begin{thm}[Achievability]\label{thm:qudit_compression_fixed_p}
Assume that the spectrum $p$ is known for the state $\rho$ to compress, and that $\rho$ has rank $r \le d$ with strictly non-degenerate non-zero eigenvalues $p_1 > p_2 > \dots > p_r > 0$. 
There exists a compression protocol of $\rho^{\otimes n}$ with memory cost
\begin{equation}
    \log|M| = \frac{r(2d-r-1)}{2}\log n + \sum_{1 \le i < j \le r}\log(p_i-p_j) + (d-r)\sum_{i=1}^r \log p_i - \sum_{k=d-r}^{d-1}\log k! + o(1)
\end{equation}
and a vanishing trace distance error $\delta=O(n^{-\frac14+\epsilon})$ for any arbitrarily small $\epsilon>0$.
\end{thm}

More explicitly, we show that there exists a compression protocol with memory cost
\begin{equation}
    \log|M| = \frac{r(2d-r-1)}{2}\log n + \sum_{1 \le i < j \le r}\log(p_i-p_j) + (d-r)\sum_{i=1}^r \log p_i - \sum_{k=d-r}^{d-1}\log k! + O\left(\frac{r d \Delta_n}{\sqrt{n}}\right)
\end{equation}
and an error 
\begin{equation}
    \delta\le O(n^{-\frac14+\epsilon}) + 2(n+1)^{\frac{r(r+1)}{2}}e^{-\Delta_n},
\end{equation}
where $\Delta_n>0$ is a slowly growing function of $n$, which we pick to be $\Delta_n=(\log n)^2$, such that the second term is superpolynomially small.

For pure states, we have
\begin{equation}\label{eq:achievable_pure}
    \log|M|=(d-1)\log n -\log (d-1)!\ .
\end{equation}

\begin{proof}
By Schur-Weyl duality, since $\rho$ has rank $r$, the state $\rho^{\otimes n}$ is strictly supported on irreducible representations with at most $r$ rows. Thus, it is unitarily equivalent to a block diagonal state
\begin{equation}
    \rho^{\otimes n}\simeq\bigoplus_{\lambda : \ell(\lambda) \le r} q_{\lambda,n}\left(\rho_\lambda\otimes\frac{I_{\map{M}_\lambda}}{m_\lambda}\right).
\end{equation}
By introducing $P_\lambda$ as the projection onto the irreducible block associated with the Young diagram $\lambda$, we have
\begin{equation}
    P_\lambda\rho^{\otimes n}P_\lambda = q_{\lambda,n}\left(\rho_\lambda\otimes\frac{I_{\map{M}_\lambda}}{m_\lambda}\right).
\end{equation}
The encoder is defined as
\begin{equation}
    \map{E}(\cdot) = \sum_{\lambda : \ell(\lambda) \le r}\map{C}_{\lambda\to \lambda^*}\tr_{\map{M}_{\lambda}}(P_\lambda(\cdot)P_\lambda),
\end{equation}
where $\map{C}_{\lambda\to\lambda^*}$ is the minimal covariant map (generalized cloner) from Proposition~\ref{prop:main}. Because $p_i = 0$ for $i > r$, we define the target representation such that it also has exactly $r$ rows:
\begin{equation}
    \lambda^*_i := \begin{cases} 
      \lceil(n+\sqrt{n}\Delta_n)p_i\rceil & \text{if } i \le r \\
      0 & \text{if } i > r 
   \end{cases}
\end{equation}
for some parameter $\Delta_n>0$ that grows with $n$. That is, we make a non-demolition measurement $\{P_\lambda\}$, discard the multiplicity part, and encode the state in the irreducible subspace of $\mathsf{U}(d)$ associated with $\lambda^*$. 

The corresponding decoder maps back to the original blocks:
\begin{equation}
    \map{D}(\cdot) = \bigoplus_{\lambda : \ell(\lambda) \le r} q_{\lambda,n}\left(\map{C}_{\lambda^*\to \lambda}(\cdot)\otimes\frac{I_{\map{M}_\lambda}}{m_\lambda}\right).
\end{equation}
The trace distance error of the compression protocol evaluates to:
\begin{equation}
\begin{aligned}
    \eps &= \frac12\|\map{D}\circ\map{E}(\rho^{\otimes n})-\rho^{\otimes n}\|_1\\
    &= \frac12\left\|\bigoplus_{\lambda}q_{\lambda,n}\left(\map{C}_{\lambda^*\to \lambda}\circ\map{C}_{\lambda\to \lambda^*}(\rho_\lambda)-\rho_\lambda\right)\otimes\frac{I_{\map{M}_\lambda}}{m_\lambda}\right\|_1\\
    &= \sum_{\lambda}q_{\lambda,n} \frac12 \left\|\map{C}_{\lambda^*\to \lambda}\circ\map{C}_{\lambda\to \lambda^*}(\rho_\lambda)-\rho_\lambda\right\|_1.
\end{aligned}
\end{equation}
Because the empirical distribution $\lambda/n$ is restricted to $r$ non-zero components, we define the typical subset of Young diagrams as:
\begin{equation}\label{eq:typical_set}
\mathcal{T}_{p,n}:=\left\{\lambda~:~\lambda_i = 0 \text{ for } i>r, \ \max_{1 \le i \le r}|\lambda_i-p_i n|\le \sqrt{n\Delta_n/2} \right\}.
\end{equation}
The error splits into typical and atypical contributions:
\begin{equation}\label{eq:error_twoterms}
\eps \le \max_{\lambda\in\mathcal{T}_{p,n}} \frac12\left\|\map{C}_{\lambda^*\to \lambda}\circ\map{C}_{\lambda\to \lambda^*}(\rho_\lambda)-\rho_\lambda\right\|_1 + \eps_{\rm tail}(q_{\lambda,n}) \ , \quad \eps_{\rm tail}(q_{\lambda,n}):=\sum_{\lambda\not\in\mathcal{T}_{p,n}}q_{\lambda,n}, 
\end{equation}
having used the maximal trace distance bound $\frac12\|\rho-\sigma\|_1\le 1$ for the atypical term.

The tail mass $\eps_{\rm tail}$ is controlled by Sanov's theorem via the following lemma \cite{christandl2006spectra,o2021quantum,yang2016efficient,hayashi2017group}:
\begin{lem}\label{lem:tail_prob}
    Let $\lambda/n$ be the empirical probability distribution associated with the Young diagram $\lambda$.
    The probability $\mathsf{Pr}[\lambda:d(\lambda/n,p)>x]$ is upper bounded by $(n+1)^{r(r+1)/2} e^{-2nx^2}$, where $d(\lambda/n,p):=(1/2)\sum_{i=1}^r|\lambda_i/n-p_i|$ is the total variation distance over the $r$-dimensional support.
\end{lem}
For any $\lambda\not\in\mathcal{T}_{p,n}$, there exists a $\lambda_i$ such that $|\lambda_i/n-p_i|>\sqrt{\Delta_n/(2n)}$. Thus, the total variation distance is strictly bounded below by $\sqrt{\Delta_n/(2n)}$, and Lemma \ref{lem:tail_prob} yields:
\begin{equation}\label{eq:tail_prob_bound}
    \eps_{\rm tail}(q_{\lambda,n})\le (n+1)^{\frac{r(r+1)}2} e^{-\Delta_n}.
\end{equation} 

For typical states $\lambda\in\mathcal{T}_{p,n}$, we bound the trace distance using the triangle inequality and the contractivity of trace distance under CPTP maps:
\begin{equation}\label{eq:error_bound1}
\begin{aligned}
    \frac12\left\|\map{C}_{\lambda^*\to \lambda}\left(\map{C}_{\lambda\to \lambda^*}(\rho_\lambda)\right)-\rho_\lambda\right\|_1 
    &\le \frac12\left\|\map{C}_{\lambda^*\to \lambda}\left(\map{C}_{\lambda\to \lambda^*}(\rho_\lambda)\right) - \map{C}_{\lambda^*\to \lambda}(\rho_{\lambda^*})\right\|_1 + \frac12\left\|\map{C}_{\lambda^*\to \lambda}(\rho_{\lambda^*}) - \rho_\lambda\right\|_1 \\
    &\le \frac12\left\|\map{C}_{\lambda\to \lambda^*}(\rho_\lambda) - \rho_{\lambda^*}\right\|_1 + \frac12\left\|\map{C}_{\lambda^*\to \lambda}(\rho_{\lambda^*}) - \rho_\lambda\right\|_1 \ .
\end{aligned}
\end{equation}

By the construction of $\lambda^*$, the ancilla size for any typical state is $\omega = \lambda^* - \lambda$. Since $\lambda_i \ge p_i n - \sqrt{n\Delta_n/2}$, the ancilla is bounded by $\|\omega\| \le \sqrt{n}\Delta_n p_1 + \sqrt{n\Delta_n/2} = O(\sqrt{n}\Delta_n)$. 
Choosing $\Delta_n = (\log n)^2$, we satisfy $\|\omega\| = O(n^{\frac12+\epsilon'})$ for any $\epsilon' > 0$. This fits exactly into the $s = 1/2$ scaling regime of Proposition~\ref{prop:main}, guaranteeing fidelities $F \ge 1 - O(n^{-\frac12 + \epsilon'})$ for both covariant maps.

Applying the Fuchs-van de Graaf inequalities, the trace distance is bounded by the fidelity as $\frac{1}{2}\|\rho-\sigma\|_1 \le \sqrt{1-F^2} \le \sqrt{2(1-F)}$. Therefore, the trace distance error for both maps is universally bounded by $\sqrt{O(n^{-\frac12+\epsilon'})} = O(n^{-\frac14+\epsilon})$. Substituting these into Eq.~(\ref{eq:error_bound1}), the error for typical states is:
\begin{equation}
    \max_{\lambda\in\mathcal{T}_{p,n}} \frac12\left\|\map{C}_{\lambda^*\to \lambda}\circ\map{C}_{\lambda\to \lambda^*}(\rho_\lambda)-\rho_\lambda\right\|_1 \le O(n^{-\frac14+\epsilon}) + O(n^{-\frac14+\epsilon}) = O(n^{-\frac14+\epsilon}) \ .
\end{equation}

Summarizing, substituting the bounds into Eq.~(\ref{eq:error_twoterms}), the total trace distance error evaluates to:
\begin{equation}
    \eps \le O(n^{-\frac14+\epsilon}) + (n+1)^{\frac{r(r+1)}2} e^{-\Delta_n} \ .
\end{equation}

Meanwhile, since the input state is mapped into the $\mathrm{GL}(d)$ irreducible subspace $\mathcal{H}_{\lambda^*}$, the memory cost is $\log|M| = \log d_{\lambda^*}$. Using the Weyl dimension formula, we split the product into blocks corresponding to the $r$ non-zero rows and the $d-r$ empty rows:
\begin{equation}
    d_{\lambda^*} = \prod_{1 \le i < j \le d} \frac{\lambda^*_i - \lambda^*_j + j - i}{j - i} = \left( \prod_{1 \le i < j \le r} \frac{\lambda^*_i - \lambda^*_j + j - i}{j - i} \right) \left( \prod_{i=1}^r \prod_{j=r+1}^d \frac{\lambda^*_i + j - i}{j - i} \right) \ .
\end{equation}
Notice that the denominator evaluates to $\prod_{1 \le i < j \le d} (j-i) = \prod_{k=1}^{d-1} k!$. The numerator of the empty-empty block ($r < i < j \le d$) evaluates exactly to $\prod_{k=1}^{d-r-1} k!$, which we have divided out. 
Bounding the terms against the ``central'' Young diagram $\lambda_0 = \lceil pn \rceil$:
\begin{equation}
    \begin{aligned}
        d_{\lambda^*} &\le \left( \prod_{1 \le i < j \le r} \frac{n(p_i - p_j)(1+\frac{\Delta_n}{\sqrt{n}}) + j - i}{j - i} \right) \left( \prod_{i=1}^r \prod_{j=r+1}^d \frac{np_i(1+\frac{\Delta_n}{\sqrt{n}}) + j - i}{j - i} \right) \\
        &= d_{\lambda_0}\left(1+O\left(\frac{rd\Delta_n}{\sqrt{n}}\right)\right).
    \end{aligned}
\end{equation}
Taking the logarithm, the leading-order memory cost is:
\begin{equation}
    \log|M| = \log d_{\lambda_0} + O\left(\frac{r d \Delta_n}{\sqrt{n}}\right) \ .
\end{equation}
Expanding $\log d_{\lambda_0}$ algebraically gives exactly $r(r-1)/2$ factors of $n$ from the first product and $r(d-r)$ factors of $n$ from the second. The total prefactor for $\log n$ is exactly $\frac{r(r-1)}{2} + r(d-r) = \frac{r(2d-r-1)}{2}$. Gathering the constants yields the final expression:
\begin{equation}
    \log|M| = \frac{r(2d-r-1)}{2}\log n + \sum_{1 \le i < j \le r}\log(p_i-p_j) + (d-r)\sum_{i=1}^r \log p_i - \sum_{k=d-r}^{d-1}\log k! + O\left(\frac{rd\Delta_n}{\sqrt{n}}\right).
\end{equation}
\end{proof}

\begin{comment}
    

\begin{thm}\label{thm:qudit_compression_fixed_p}
Assume that the (non-degenerate) spectrum $p$ is known for the state $\rho$ to compress. 
There exists a compression protocol of $\rho^{\otimes n}$ with memory cost
\begin{equation}
    \log|M| = \frac{d(d-1)}{2}\log n+\sum_{i<j}\log(p_i-p_j) -\sum_{k=1}^{d-1}\log k! + o\left(1\right)
\end{equation}
and a vanishing error $\delta=O(n^{-\frac14+\epsilon})$ with any $\epsilon>0$.
\end{thm}
More explicitly, we show that there exists a compression protocol with memory cost
\begin{equation}
    \log|M| = \frac{d(d-1)}{2}\log n+\sum_{i<j}\log(p_i-p_j) -\sum_{k=1}^{d-1}\log k! + O\left(\frac{d^2\Delta_n}{\sqrt{n}}\right)
\end{equation}
and an error 
\begin{equation}
    \delta\le O(n^{-\frac12+\epsilon}) + 2(n+1)^{\frac{d(d+1)}{2}}e^{-\Delta_n},
\end{equation}
where $\Delta_n>0$ is a slowly growing function of $n$, which we pick to be $\Delta_n=(\log n)^2$, such that the second term is superpolynomially small.
\begin{proof}
To begin with, we notice that $\rho^{\otimes n}$ is unitarily equivalent to a block diagonal state
\begin{equation}
    \rho^{\otimes n}\simeq\bigoplus_{\lambda}q_{\lambda,n}\left(\rho_\lambda\otimes\frac{I_{\map{M}_\lambda}}{m_\lambda}\right).
\end{equation}
By introducing $P_\lambda$ as the projection on the irreducible block associated with the Young diagram $\lambda$, we have
\begin{equation}
    P_\lambda\rho^{\otimes n}P_\lambda = q_{\lambda,n}\left(\rho_\lambda\otimes\frac{I_{\map{M}_\lambda}}{m_\lambda}\right).
\end{equation}
For some (vanishing) parameter $\xi>0$, the encoder is defined as
\begin{equation}
    \map{E}(\cdot) = \sum_{\lambda}\map{C}_{\lambda\to \lambda^*}\tr_{\map{M}_{\lambda}}(P_\lambda(\cdot)P_\lambda),
\end{equation}
where $\map{C}_{\lambda\to\lambda^*}$ is the generalised universal quantum cloner \eqref{eq:gen_cloner} and
\begin{equation}
    \lambda^*:=\lceil(n+\sqrt{n}\Delta_n)p\rceil
\end{equation}
for some parameter $\Delta_n>0$ that grows with $n$.
That is, for encoding, we first make a non-demolition measurement $\{P_\lambda\}$ and discard the multiplicity part that contains no information about $\rho$. We then proceed with the universal cloner to encode the state in the irreducible subspace of $\mathsf{U}(d)$ associated with $\lambda^*$. 

The corresponding decoder is defined as
\begin{equation}
    \map{D}(\cdot) = \bigoplus_{\lambda}q_{\lambda,n}\left(\map{C}_{\lambda^*\to \lambda}(\cdot)\otimes\frac{I_{\map{M}_\lambda}}{m_\lambda}\right).
\end{equation}
The error of the compression protocol for an arbitrary input $\rho^{\otimes n}$ can be expressed as
\begin{equation}
\begin{aligned}
    \eps=&\frac12\|\map{D}\circ\map{E}(\rho^{\otimes n})-\rho^{\otimes n}\|_1\\
    =&\frac12\left\|\bigoplus_{\lambda}q_{\lambda,n}\left(\map{C}_{\lambda^*\to \lambda}\circ\map{E}(\rho^{\otimes n})-\rho_\lambda\right)\otimes\frac{I_{\map{M}_\lambda}}{m_\lambda}\right\|_1\\
     =&\frac12\sum_{\lambda}q_{\lambda,n}\left\|\map{C}_{\lambda^*\to \lambda}\circ\map{E}(\rho^{\otimes n})-\rho_\lambda\right\|_1.
\end{aligned}
\end{equation}
Now, we define a typical subset of the Young diagrams with $n$ boxes and $d$ rows as:
\begin{equation}\label{eq:typical_set}
\mathcal{T}_{p,n}:=\left\{\lambda~:~\max_{i}|\lambda_i-p_i n|\le \sqrt{n\Delta_n/2} \right\}.
\end{equation}
The error can be split into two terms and bounded as
\begin{equation}\label{eq:error_twoterms}
\eps\le\max_{\lambda\in\mathcal{T}_{p,n}}\frac12\left\|\map{C}_{\lambda^*\to \lambda}\circ\map{E}(\rho^{\otimes n})-\rho_\lambda\right\|_1+\eps_{\rm tail}(q_{\lambda,n})\qquad \eps_{\rm tail}(q_{\lambda,n}):=\sum_{\lambda\not\in\mathcal{T}_{p,n}}q_{\lambda,n}, 
\end{equation}
having used the bound $(1/2)\|\rho-\sigma\|_1\le 1$ for two arbitrary states $\rho,\sigma$.

We can bound both terms separately. 
The second term can be bounded by applying the following lemma \cite{christandl2006spectra,o2021quantum,yang2016efficient,hayashi2017group}:
\begin{lem}\label{lem:tail_prob}
    Let $\lambda/n$ be the probability distribution $(\lambda_i/n)$ associated with the Young diagram $\lambda$ and $p$ be the spectrum of $\rho$.
    The probability 
    $$\mathsf{Pr}[\lambda:d(\lambda/n,p)>x]:=\sum_{\lambda:d(\lambda/n,p)>x}q_{\lambda,n}$$ is upper bounded by $(n+1)^{d(d+1)/2}\cdot e^{-2nx^2}$, where $d(\lambda/n,p):=(1/2)\sum_i|\lambda_i/n-p_i|$ is the total variation distance.
\end{lem}
By the definition (\ref{eq:typical_set}) of the typical set $\mathcal{T}_{p,n}$, for any $\lambda\not\in\mathcal{T}_{p,n}$ there exists a $\lambda_i$ such that $|\lambda_i/n-p_i|>\sqrt{\Delta_n/(2n)}$ and thus the total variation distance is higher than $\sqrt{\Delta_n/(2n)}$. Therefore, Lemma \ref{lem:tail_prob} yields
\begin{equation}\label{eq:tail_prob_bound}
    \eps_{\rm tail}(q_{\lambda,n})\le (n+1)^{\frac{d(d+1)}2}\cdot e^{-\Delta_n}.
\end{equation} 
The first term in Eq.~(\ref{eq:error_twoterms}) can be expanded and bounded as
\begin{equation}\label{eq:error_bound1}
\begin{aligned}
    \max_{\lambda\in\mathcal{T}_{p,n}}\frac12\left\|\map{C}_{\lambda^*\to \lambda}\circ\map{E}(\rho^{\otimes n})-\rho_\lambda\right\|_1& = \max_{\lambda\in\mathcal{T}_{p,n}}\frac12\left\|\sum_{\lambda'}q_{\lambda',n}\map{C}_{\lambda^*\to \lambda}\circ\map{C}_{\lambda'\to \lambda^*}(\rho_{\lambda'})-\rho_\lambda\right\|_1\\
    &\le \max_{\lambda\in\mathcal{T}_{p,n}}\frac12\sum_{\lambda'}q_{\lambda',n}\left\|\map{C}_{\lambda^*\to \lambda}\circ\map{C}_{\lambda'\to \lambda^*}(\rho_{\lambda'})-\rho_\lambda\right\|_1\\
    &\le \max_{\lambda',\lambda\in\mathcal{T}_{p,n}}\frac12\left\|\map{C}_{\lambda^*\to \lambda}\circ\map{C}_{\lambda'\to \lambda^*}(\rho_{\lambda'})-\rho_\lambda\right\|_1+\eps_{\rm tail}(q_{\lambda,n}).
\end{aligned}
\end{equation}
The first inequality comes from the convexity of the trace distance (i.e. $\|\sum_\lambda q_\lambda(\cdot)\|_1\le \sum_\lambda q_\lambda\|(\cdot)\|_1$) and the second inequality can be derived in the same way as Eq.~(\ref{eq:error_twoterms}).


For $\lambda\in\mathcal{T}_{p,n}$, $\lambda^*-\lambda$ is a Young diagram of size $O(\sqrt{n}\Delta_n)$, as long as $\Delta_n>\max_{i,j}2/(p_i-p_j)^2$.
This satisfies the condition of Proposition~\ref{prop:main} when $\Delta_n\ll \sqrt{n}$, and $\map{C}_{\lambda\to\lambda^*}$ can faithfully convert $\rho_\lambda$ and $\rho_{\lambda^*}$ (and vice versa). 

Explicitly, for $\Delta_n=(\log n)^2$, we have $\|\omega\|=O(n^{\frac12+\epsilon})$ for any $\epsilon>0$, and the fidelity in Proposition~\ref{prop:main} is lower bounded by $1-O(n^{-\frac12+\epsilon})$. Therefore, we get
\begin{equation}
\begin{aligned}
&\max_{\lambda',\lambda\in\mathcal{T}_{p,n}}\frac12\left\|\map{C}_{\lambda^*\to \lambda}\circ\map{C}_{\lambda'\to \lambda^*}(\rho_{\lambda'})-\rho_\lambda\right\|_1\\
\le &\max_{\lambda\in\mathcal{T}_{p,n}}\frac12\left\|\map{C}_{\lambda^*\to \lambda}(\rho_{\lambda^*})-\rho_\lambda\right\|_1+\max_{\lambda',\lambda\in\mathcal{T}_{p,n}}\frac12\left\|\map{C}_{\lambda^*\to \lambda}\circ\map{C}_{\lambda'\to \lambda^*}(\rho_{\lambda'})-\map{C}_{\lambda^*\to \lambda}(\rho_{\lambda^*})\right\|_1\\ 
\le &\max_{\lambda\in\mathcal{T}_{p,n}}\frac12\left\|\map{C}_{\lambda^*\to \lambda}(\rho_{\lambda^*})-\rho_\lambda\right\|_1+\max_{\lambda'\in\mathcal{T}_{p,n}}\frac12\left\|\map{C}_{\lambda'\to \lambda^*}(\rho_{\lambda'})-\rho_{\lambda^*}\right\|_1  \le O(n^{-\frac14+\epsilon})
\end{aligned}
\end{equation}
Summarizing, substituting Eqs.~(\ref{eq:tail_prob_bound}) and (\ref{eq:error_bound1}) into Eq.~(\ref{eq:error_twoterms}), we obtain 
\begin{equation}
\begin{aligned}
    \eps&\le\max_{\lambda',\lambda\in\mathcal{T}_{p,n}}\frac12\left\|\map{C}_{\lambda^*\to \lambda}\circ\map{C}_{\lambda'\to \lambda^*}(\rho_{\lambda'})-\rho_\lambda\right\|_1+2\eps_{\rm tail}(q_{\lambda,n})\\
    &\le O(n^{-\frac14+\epsilon}) +2 (n+1)^{\frac{d(d+1)}2}\cdot e^{-\Delta_n}.
\end{aligned}
\end{equation}

Meanwhile, since the input state is mapped into $\mathcal{U}_{\lambda^*}$, the cost of the compression is
\begin{equation}
    \log|M|=\log d_{\lambda^*}.
\end{equation}
The dimension of this irreducible subspace can be bounded as
\begin{equation}
    \begin{aligned}
        d_{\lambda^*}&=\frac{\prod_{i<j}(\lambda^*_i-\lambda^*_j-i+j)}{\prod_{k=1}^{d-1}k!}\\
        &\le\frac{\prod_{i<j}((1+\Delta_n/\sqrt{n})n(p_i-p_j)+1-i+j)}{\prod_{k=1}^{d-1}k!}\\
        &= d_{\lambda_0}(1+O(d^2\Delta_n/\sqrt{n})).
    \end{aligned}
\end{equation}
Here $\lambda_0=\lceil pn\rceil$ denotes the ``central" Young diagram defined by the spectrum $p$. Therefore, the memory cost is
\begin{equation}
    \log|M|=\log d_{\lambda_0}+O\left(\frac{d^2\Delta_n}{\sqrt{n}}\right).
\end{equation}
\end{proof}

\end{comment}


\section{Converse}
\begin{comment}
    

\subsection{Review of YCH on the converse from Holevo information}
YCH argued the converse through the Holevo information $\chi_H$. Consider the continuous Haar ensemble $E:=\{g,(U(g)\rho U(g)^\dagger)^{\otimes n}\}$, and and compression code $(\E,\D)$. The data processing inequality implies
\begin{equation}
    \chi_H(E)\ge \chi_H(\E(E):=\{g,\E(U(g)\rho U(g)^\dagger)^{\otimes n})\})\ge \chi_H(\D\circ\E(E):=\{g,\D\circ\E(U(g)\rho U(g)^\dagger)^{\otimes n})\})
\end{equation}
The criteria~\eqref{eq:error} implies that 
\begin{equation}\label{eq:continuity}
    \chi_H(\D\circ\E(E))\ge \chi_H(E) - 2\eps\log d_E +2\eps\log\eps
\end{equation}
where the effective dimension $d_E$ is the rank of the ensemble mixture. It is exponential in $n$, so instead YCH argued that $d_E$ can be replaced by $d_E^{\min}$, which is the minimum of $d_{E'}$ over all ensembles $E'$ that are sufficient statistics for $E$. (We say $E'$ is a sufficient statistic for $E$ if $E$ can be encoded into $E'$ and decoded with zero error, and $\chi(E)=\chi(E')$.) The upshot is that $d_E^{\min}$ can be polynomial in $n$.

We thus have the lower bound on the memory size
\begin{equation}\label{eq:chain}
    \log |M|\ge \chi_H(\E(E))\ge \chi_H(E) - 2\eps\log d^{\min}_E +2\eps\log\eps=\chi_H(E)-\O(\eps\log n)+\O(\eps)\ .
\end{equation}

Now it remains to estimate the Holevo information for the Haar ensemble $E$. This quantity can be computed as follows. The ensemble mixture $\int \dd g\, (U(g)\rho U(g)^\dagger)^{\otimes n}$ is invariant under both the unitary group and the permutation group. Schur's lemma implies that it is flat within each block, so it is unitarily equivalent to
\begin{equation}
    \bigoplus_\lambda p_\lambda\, \tau_{V_\lambda}\otimes\tau_{\mathcal{V}_\lambda},\quad p_\lambda = \chi_\lambda(\rho)\cdot\dim\mathcal{V}_\lambda\ ,
\end{equation}




Then we have
\begin{multline}
   \chi_H(E) =  \sum_\lambda p_\lambda(\log\dimV_\lambda+\log\dim\mathcal{V}_\lambda-\log p_\lambda)  -nH(\rho)\\
   =\sum_\lambda p_\lambda(\log\dimV_\lambda-\log\chi_\lambda(\rho)) -nH(\rho)=\sum_\lambda p_\lambda\log\dimV_\lambda -f(\mathbf{p})\ .
\end{multline}
In the last step, the $n$-dependent terms cancel so $\sum_\lambda p_\lambda\log\chi_\lambda(\rho)+nH(\rho)=f(\mathbf{p})\sim \O(1)$, where~$\mathbf{p}$ denotes the spectrum of $\rho$. 

Alternatively, we can use a sufficient statistic of $E$ to obtain the same result. Consider again the Schur-Weyl decomposition~\eqref{eq:schur_d}. We now trace out the multiplicity register.
\begin{equation}
    V(U(g)\rho U(g)^\dagger)^{\otimes n} V^\dagger = \bigoplus_\lambda p_\lambda \rho_{g,\lambda}\ .
\end{equation}
Consider now the ensemble, $E':=\{g, \bigoplus_\lambda p_\lambda \rho_{g,\lambda}\}$ which is a sufficient statistic of $E$, so $\chi_H(E')=\chi_H(E)$. The ensemble mixture is again flat within each block. Let us calculate $\chi_H(E')$,
\begin{equation}
    \chi_H(E') = H\left(\bigoplus_\lambda p_\lambda \tau_{V_\lambda}\right) - H\left(\bigoplus_\lambda p_\lambda \rho_\lambda\right) = \sum_\lambda p_\lambda\log\dimV_\lambda - \sum_\lambda p_\lambda H(\rho_\lambda)\ .
\end{equation}
It follows that 
\begin{equation}\label{eq:remainder}
    \sum_\lambda p_\lambda H(\rho_\lambda)=\sum_\lambda p_\lambda\log\chi_\lambda(\rho)+nH(\rho)=f(\mathbf{p})\sim \O(1)\ .
\end{equation}


We now have the converse bound
\begin{equation}\label{eq:converse}
    \log |M|\ge \left(\frac{d(d-1)}{2}-\eps\right)\log n+\sum_{i<j}\log(p_i-p_j) -\sum_{k=1}^{d-1}\log k!-f(\mathbf{p})+\O(\eps)\ .
\end{equation}
The converse matches the achievability upper bound~\eqref{eq:achievable} up to the $f(\mathbf{p})>0$ remainder term. This remainder term makes the converse a bit looser. It's plausible that for a fixed eigenvalue density $\mu$ to which $\mathbf{p}$ asymptotes to at large $d$, $f(\mathbf{p})$ doesn't scale with $d$ so that this offset becomes negligible at large $d$. For pure states, this term vanishes and the converse exactly matches the achievability.





Note that \eqref{eq:chain} can only be tight if the codewords in $\E(E)$ are pure. Conversely, \eqref{eq:converse} cannot never tight if the codewords have some $\O(1)$ entropy, and the remainder term is necessary. We can tighten the lower bound by adding to the RHS of \eqref{eq:converse} the average entropy of the codewords. This extra term, however, depends on the encoder $\E$ under consideration. It is tricky to argue for a universal lower bound for the entropy of codewords for any encoder. 

Nonetheless, it is instructive to see if the YCH encoder itself actually gives a tight lower bound. At large $n$, the probability distribution $p_\lambda$ concentrates, so 
\begin{equation}
     f(\mathbf{p})=\sum_\lambda p_\lambda H(\rho_\lambda)\to H(\rho_{\lambda^*})\ .
\end{equation}
The average entropy of the codewords should be well approximated by $H(\rho_{\lambda^*})$, so the lower bound for the specific encoder can be tightened to match~\eqref{eq:achievable_qubit}. \\

\end{comment}



\begin{thm}[Converse]
    Any quantum compression protocol that asymptotically compresses the ensemble $\mathcal{E} = \{dg, \rho_g^{\otimes n}\}$ with vanishing error must use a memory $M$ whose size is asymptotically lower bounded by:
    \begin{equation}
        \log|M| \ge \frac{r(2d-r-1)}{2}\log n + \sum_{1 \le i < j \le r}\log(p_i-p_j) + (d-r)\sum_{i=1}^r \log p_i - \sum_{k=d-r}^{d-1}\log k! - o(1) \ .
    \end{equation}
    This provides a tight converse that strictly matches the achievability upper bound.
\end{thm}

\begin{proof}
Consider the Schur-Weyl block diagonalization of the input states. The completely mixed ensemble is unitarily equivalent to:
\begin{equation}
    V \left( \int \mathrm{d}g \, U(g)\rho^{\otimes n} U(g)^\dagger \right) V^\dagger = \bigoplus_{\lambda} q_{\lambda,n} \left( \rho_{g,\lambda} \otimes \tau_\lambda \right) \ ,
\end{equation}
where $\rho_{g,\lambda} = \pi_\lambda(g) \rho_{e,\lambda} \pi_\lambda(g)^\dagger \in \mathcal{D}(V_\lambda)$ acts on the irreducible representation space, and $\tau_\lambda \in \mathcal{D}(\mathcal{M}_\lambda)$ acts on the multiplicity space.

To bound the memory size, we consider a $\Lambda$-assisted protocol. Suppose we grant the decoder free access to a classical register $\Lambda$ containing the exact value of the measurement outcome $\lambda$ (which occurs with probability $q_{\lambda,n}$). Because providing free classical side-information to the decoder can only reduce or maintain the required quantum memory size, any lower bound on the memory of this $\Lambda$-assisted task constitutes a strict lower bound on the memory of the original unassisted task.

Conditioned on receiving a specific $\lambda$, the task reduces to compressing the mixed state ensemble $\mathcal{E}_\lambda = \{dg, \rho_{g,\lambda}\}$. According to the Koashi-Imoto compression theorem \cite{koashi2001compressibility,koashi2001possible}, the optimal asymptotic memory size for a mixed-state ensemble is dictated by the dimension of its ``pruned'' or ``core'' system. The Koashi-Imoto decomposition splits the Hilbert space into a direct sum of tensor products $\bigoplus_j \mathcal{H}_j^L \otimes \mathcal{H}_j^R$, where the states in the ensemble vary only over the left systems $\mathcal{H}_j^L$, leaving the right systems $\mathcal{H}_j^R$ as state-independent redundant noise that can be discarded.

For the ensemble $\mathcal{E}_\lambda$, the states are generated by the continuous group action $\pi_\lambda(g)$. The $C^*$-algebra $\mathcal{A}_\lambda$ generated by the support of the states $\{\rho_{g,\lambda}\}_{g \in \mathrm{U}(d)}$ is therefore invariant under the unitary representation $\pi_\lambda(g)$. 

Crucially, because $\pi_\lambda$ is an \emph{irreducible} representation of $\mathrm{U}(d)$, Schur's Lemma dictates that it possesses no non-trivial invariant subspaces. Consequently, the invariant algebra $\mathcal{A}_\lambda$ must be the entire bounded operator algebra $\mathcal{B}(V_\lambda)$. Because the algebra spans the entire space, the Koashi-Imoto decomposition is completely trivial: there are no redundant $R$-registers or classical $j$-blocks to discard. The pruned core is exactly the full irreducible space~$V_\lambda$.

Therefore, conditioned on $\lambda$, the optimal asymptotic memory size required to compress $\mathcal{E}_\lambda$ with vanishing error is exactly $\log \dim V_\lambda$.

In the standard formulation of quantum compression, the protocol must use a single, fixed-length memory $|M|$ that succeeds for the entire ensemble with high probability. By standard typicality arguments, the probability distribution $q_{\lambda,n}$ is highly concentrated around the typical representations $\mathcal{T}_{p,n}$ where $\lambda_i = n p_i \pm O(\sqrt{n})$ for $1 \le i \le r$, and strictly $\lambda_i = 0$ for $r < i \le d$. 

To guarantee vanishing error, the fixed memory capacity must be large enough to accommodate the maximum dimension within this typical set. To evaluate the explicit asymptotic scaling, we expand the Weyl dimension formula evaluated at these typical representations:
\begin{equation}
    \dim V_\lambda = \frac{\prod_{1 \le i < j \le d} (\lambda_i - \lambda_j + j - i)}{\prod_{1 \le i < j \le d} (j - i)} \ .
\end{equation}
We split the product in the numerator into three distinct blocks based on the effective rank~$r$:
\begin{enumerate}
    \item \textbf{Non-zero vs. Non-zero rows ($1 \le i < j \le r$):} 
    The terms scale as $\lambda_i - \lambda_j = n(p_i - p_j) \pm O(\sqrt{n})$. Factoring out $n(p_i - p_j)$, the contribution to the product is:
    \begin{equation}
        \prod_{1 \le i < j \le r} n(p_i - p_j) \left( 1 + O\left(\frac{1}{\sqrt{n}}\right) \right) = n^{\frac{r(r-1)}{2}} \left( \prod_{1 \le i < j \le r} (p_i - p_j) \right) \left( 1 + O\left(\frac{1}{\sqrt{n}}\right) \right) \ .
    \end{equation}
    
    \item \textbf{Non-zero vs. Empty rows ($1 \le i \le r < j \le d$):}
    Since $\lambda_j = 0$, the terms scale as $\lambda_i + j - i = n p_i \pm O(\sqrt{n})$. Factoring out $n p_i$, and noting there are exactly $r(d-r)$ such terms, the contribution is:
    \begin{equation}
        \prod_{i=1}^r \prod_{j=r+1}^d n p_i \left( 1 + O\left(\frac{1}{\sqrt{n}}\right) \right) = n^{r(d-r)} \left( \prod_{i=1}^r p_i \right)^{d-r} \left( 1 + O\left(\frac{1}{\sqrt{n}}\right) \right) \ .
    \end{equation}
    
    \item \textbf{Empty vs. Empty rows ($r < i < j \le d$):}
    Both row lengths are strictly zero, so the terms are exactly $j - i$. This product evaluates to a product of factorials:
    \begin{equation}
        \prod_{r < i < j \le d} (j - i) = \prod_{k=1}^{d-r-1} k! \ .
    \end{equation}
\end{enumerate}

The denominator of the Weyl dimension formula is exactly $\prod_{k=1}^{d-1} k!$. Dividing the exact factorial block from the numerator (Block 3) by the denominator leaves:
\begin{equation}
    \frac{\prod_{k=1}^{d-r-1} k!}{\prod_{k=1}^{d-1} k!} = \frac{1}{\prod_{k=d-r}^{d-1} k!} \ .
\end{equation}

Multiplying the polynomial factors of $n$ from Blocks 1 and 2 gives $n^{\frac{r(r-1)}{2}} n^{r(d-r)} = n^{\frac{r(2d-r-1)}{2}}$. Combining all pieces, the typical dimension scales as:
\begin{equation}
    \dim V_\lambda = n^{\frac{r(2d-r-1)}{2}} \frac{\prod_{i < j \le r}(p_i - p_j) \cdot \prod_{i=1}^r p_i^{d-r}}{\prod_{k=d-r}^{d-1} k!} \left( 1 + O\left(\frac{1}{\sqrt{n}}\right) \right) \ .
\end{equation}
Taking the logarithm yields the asymptotic memory cost. The $O(1/\sqrt{n})$ fluctuations inside the logarithm become $\log(1 \pm O(1/\sqrt{n})) = \pm O(1/\sqrt{n}) = o(1)$:
\begin{equation}
    \log \dim V_\lambda = \frac{r(2d-r-1)}{2}\log n + \sum_{1 \le i < j \le r}\log(p_i-p_j) + (d-r)\sum_{i=1}^r \log p_i - \sum_{k=d-r}^{d-1}\log k! \pm o(1) \ .
\end{equation}
Because the unassisted protocol's required memory $|M|$ must be at least the expected value over the typical set, we have $\log|M| \ge \sum_\lambda q_{\lambda,n} \log \dim V_\lambda - o(1)$. This establishes the tight converse bound matching the achievability result.
\end{proof}

\section{The visible setting: Programming quantum states}
We could consider a variant of the compression task, where the state to be compressed is \emph{visible} to the encoder. This is in contrast to the \emph{blind} setting that we have been studying. Now, we no longer restrict the encoder to use one universal channel $\E$. Instead, she can use a different encoding for each state to be compressed. 

Equivalently, this setting can be reformulated as \emph{programming} quantum states. The task is to find the minimal quantum memory cost for programming any instance of $\{\rho^{\otimes n}_U\}_U$, such that a universal state preparation device, modeled as some CPTP map $\D$, can output the correct state,
\begin{equation}
    \D(\sigma_{\rho^{\otimes n}_U}) = \rho^{\otimes n}_U\ ,\quad\forall\, U\in\mathrm{U}(d)
\end{equation}
where $\sigma_{\rho^{\otimes n}_U}$ denotes the program for $\rho^{\otimes n}_U$.

The achievability part should be essentially the same as in the blind setting.

The converse part is more difficult. \JW{Idea: We first show that $\D$ can be assumed covariant w.l.o.g.; and then the converse follows from the rigidity imposed by Schur's lemma.}

\section{Free entropy of unitary operators: Unitary programming}

One key distinction between quantum entropy theory and free entropy theory is that the former is based on generalizing probability distribution density to the non-commutative setting, whereas free entropy is based on generalizing probability distribution itself to the non-commutative setting. Therefore, unlike the von Neumann entropy, free entropy is also defined for operators other than density operators, such as unitary operators and Hermitian operators of non-unit trace.

We ask for a programming scheme consisting of a program (quantum state) and a universal processor (CPTP map), $(\sigma_U,\mathcal{P})$, that can \emph{simulate} the programmed unitary $U$,
\begin{equation}
    ||\mathcal{P}(\,\cdot\,\otimes\sigma_U) - U\,\cdot\,U^\dagger||_\diamond\le\eps
\end{equation}

This problem has been studied by Yang et al~\cite{yang2020optimal} for the full set of unitary operations in finite dimensions. Here we need to slightly generalize their result to subsets of unitaries parameterized by fewer parameters. This could help us deal with situations where the spectrum of the unitary is known and possibly degenerate.

We will show that the optimal unitary program has a size determined by the physical free entropy up to the subleading $\O(1)$ term. \JW{Any idea on how to obtain the $\O(1)$ term? Kirillov character formula?}

\subsection{Free entropy of unitary}
We can define the free entropy of a unitary operator by measuring the volume of unitary matricial microstates (cf. Hiai-Petz \cite{hiai2000semicircle}. The resulting formula is essentially the same up to changing the integration domain from $\mathbb{R}$ to $\mathbb{T}$. 
\begin{equation}
    \chi(u)=\int_\mathbb{T}\int_\mathbb{T}\log|z-z'|\dd\mu_u(z)\dd\mu_u(z')+ \mathrm{const}
\end{equation}
Similarly, we can define the physical free entropy as,
\begin{equation}
    \chi_\mathrm{phy}(u;\epsilon):=\delta(u)\log\epsilon^{-1} + \chi_\mathrm{reg}(u)
\end{equation}
where $u$ is the free entropy dimension of $u$, and $\chi_\mathrm{reg}(u)$ is the finite part of $\chi(u)$. Explicitly, for a $d$-dimensional unitary operator,
\begin{equation}
    \chi_\mathrm{phy}(u;\epsilon) = \left(1-\sum_i \frac{g_i^2}{d^2}\right)\log\epsilon^{-1} + \frac{2}{d^2}\sum_{i < j}g_ig_j\log|z_i-z_j|
\end{equation}
where $i$ labels distinct eigenvalues and $g_i$'s denote the eigenvalue degeneracies.




\JW{@Yuxiang, could you fill these sections below?}
\subsection{Achievability}\label{subsec:programmingunitary}
Consider an arbitrary $f$-parameter unitary family $\mathbb{F}$. That is, there exists an one-to-one and onto function from $\mathbb{T}^f$ to $\mathbb{F}$.  In the following, we show an explicit programming scheme of unitary in $\mathbb{F}$ with error $\epsilon$ and cost $(f/2)\log\epsilon^{-1}$ in the leading order.
 
To show this, we need to come up with an explicit programming scheme, which runs as follows:
\begin{itemize}
    \item For any gate $U_x\in\mathbb{F}$ with $x=(x_1,x_2,\dots,x_f)\in\mathbb{T}^f$, we could use $f$ independent sine-shape states \cite{buvzek1999optimal} as the program state:
    \begin{align*}
        \ket{P_x}=\ket{\sin_{x_1}}\otimes\cdots\otimes\ket{\sin_{x_f}}
    \end{align*}
    with $\ket{\sin_{y}}:=\sqrt{2/D}\sum_{y=0}^{D-1} \sin\frac{(m+1/2)\pi}{D}e^{imy}\ket{m}$ and $D$ being a suitable parameter to be determined later, each to encode one individual degree of freedom. 
    \item To execute the unitary, we just use a measure-and-operate strategy to implement the corresponding gate, i.e., we measure each $\ket{\sin_{x_i}}$ in the Fourier basis of $D$ dimension and operate the unitary according to the estimate.  
\end{itemize}
Regarding the performance, the programming cost is $f\log D+O(1)$. For the error, since the sine-shape state achieves the Heisenberg limit with respect to $D$ \cite{buvzek1999optimal}, the estimation error will scale at most as $1/D^2$ in terms of MSE (mean squared error). Details are provided in the following subsection. Combining these arguments, we get $(f/2)\log\eps^{-1}$ as the leading order of the programming cost.

\subsection{Error of estimation-based programming}
For any $x$ and small $y$, we have (\YY{Here we need to assume some smoothness and/or continuity}):
\begin{align}\label{eq:taylorUx}
    \mathcal{U}_{x+y} = \mathcal{U}_{x} + \sum_{k=1}^{\infty}\sum_{j_1,\dots,j_k=1}^{f}y_{j_1}\cdots y_{j_k}\frac{\partial^k \mathcal{U}_x}{\partial x_{j_1}\cdots\partial x_{j_k}},
\end{align}
where all derivatives are defined at $x$. We further assume that the $k$-th order derivative is bounded as
\begin{align}\label{eq:derivativebound}
    \left\|\frac{\partial^k \mathcal{U}_x}{\partial x_{j_1}\cdots\partial x_{j_k}}\right\|_{\diamond}\le (2h)^k
\end{align}
for any $k$, where $h>0$ is a universal constant. In the case where $U_x=\exp\{-iH\cdot x\}$ for some $x$-independent Hamiltonians $H:=(H_1,\dots,H_f)^T$, one has $\partial \mathcal{U}_x/\partial{x_j}(\cdot)=-i[H_j,\partial \mathcal{U}_x(\cdot)]$, and it is sufficient to set $h=\max_j\|H_j\|$.

Any estimation protocol of $x\in\mathbb{T}^f$ is characterized by a series of conditional probability distributions $d\hat{x}\,p(\hat{x}|x)$, which is the distribution of the estimate $\hat{x}$ when the true parameter is $x$. The error of the estimation protocol is characterized by the worst-case mean-squared error (MSE):
\begin{align}
    \epsilon_{\rm est}:=\sup_{x}\int d\hat{x}\,p(\hat{x}|x)\sum_{j=1}^f(\hat{x}_j-x_j)^2.
\end{align}

Our bound on the error of the specific programming protocol based on the sine state is a special case of the following more general result, which sets a sufficient condition for the equivalence of asymptotic programming and unitary estimation:
\begin{lem}\label{lem:estprog}
    Consider any estimation protocol of $x$ with error $\epsilon\ll 1$ and a conditional probability density function $p(\hat{x}|x)$ satisfying the following conditions: 
    \begin{enumerate}
        \item there exists a probability density function $q$ such that $p(\hat{x}|x)=q(\hat{x}-x)$ for any $\hat{x},x$, 
        \item $\exists\,\alpha>0$ such that $1-\int_{T_\epsilon}dy~q(y)=o(\epsilon)$ for $T_\epsilon:=\{y:y_j\in[-\epsilon^\alpha,\epsilon^\alpha],\forall\,j\}$,
        \item  $q(y)$ is an even function within $T_\eps$, i.e., $q(y)=q(-y)$ for any $y\in T_\eps$.
    \end{enumerate}
    Then, the estimation-based programming protocol $\int d\hat{x}\,p(\hat{x}|x)\,\mathcal{U}_{\hat{x}}$ has error
    \begin{align}
        \eps_{\rm prog}\le \epsilon\cdot\frac12\left\|\frac{\partial^2\map{U}_x}{\partial x_{j}^2}\right\|_{\diamond} + o(\epsilon).
    \end{align}
\end{lem}
\begin{proof} 
    The programming error is expressed as:
    \begin{align}
        \eps_{\rm prog}:=\sup_{x}\frac12\left\|\int d\hat{x}\,p(\hat{x}|x)\mathcal{U}_{\hat{x}}-\mathcal{U}_x\right\|_{\diamond}.
    \end{align}
    By switching between parametrizations (condition 1) and the vanishing tail condition (condition 2), we can bound it as:
    \begin{align}
        \eps_{\rm prog}\le\sup_{x}\frac12\left\|\int_{T_\eps} dy\,q(y)\left(\mathcal{U}_{x+y}-\mathcal{U}_x\right)\right\|_{\diamond}+o(\eps).
    \end{align}
    Since $T_\epsilon$ is a small neighborhood around $x$, we can expand $\mathcal{U}_{x+y}$ into a Taylor series.
    Substituting Eq.~(\ref{eq:taylorUx}) into the above bound, we get
     \begin{align}
        \eps_{\rm prog}\le\sup_{x}\frac12\left\|\sum_{k=1}^{\infty}\sum_{j_1,\dots,j_k=1}^{f}\left(\int_{T_\eps} dy\,q(y)y_{j_1}\cdots y_{j_k}\right)\frac{\partial^k \mathcal{U}_x}{\partial x_{j_1}\cdots\partial x_{j_k}}\right\|_{\diamond}+o(\eps).
    \end{align}
    Recall that for a group of real numbers $\{c_j\}$ and linear maps $\{\mathcal{A}_j\|_{\diamond}$, $\|\sum_j c_j\mathcal{A}_j\|_{\diamond}\le\sum_j|c_j|\|\mathcal{A}_j\|_{\diamond}$, which can be shown via the triangle inequality and the absolute homogeneity of the diamond norm. Therefore, the programming error can be further bounded as:
    \begin{align}\label{eq:progerr1}
        \eps_{\rm prog}\le\sup_{x}\sum_{k=1}^{\infty}\sum_{j_1,\dots,j_k=1}^{f}\left|\int_{T_\eps} dy\,q(y)y_{j_1}\cdots y_{j_k}\right|\frac12\left\|\frac{\partial^k \mathcal{U}_x}{\partial x_{j_1}\cdots\partial x_{j_k}}\right\|_{\diamond}+o(\eps).
    \end{align}
    
    Notice that $T_\eps$ is symmetric by definition. Therefore, since $q(y)$ is an even function within $T_\eps$ (condition 3), we have
    \begin{align}
        \int_{T_\eps}dy\,q(y)y_{j}=0
    \end{align}
    for any $j$. That is, the estimation protocol is locally unbiased in $T_\eps$. 
    More generally, consider arbitrary fixed $k$. When $k$ is odd, we have 
    \begin{align}
        \int_{T_\eps}dy\,q(y)(y_{j})^k=0;
    \end{align}
    when $k$ is even, we can bound the $k$-moment as
    \begin{align}
        \int_{T_\eps}dy\,q(y)(y_{j})^k\le \max_{y\in T_\eps}|y_j|^{k-2} \int_{T_\eps}dy\,q(y)(y_{j})^2 = \epsilon^{1+(k-2)\alpha}.
    \end{align}
    With these, further consider arbitrary indices $j_1,\dots,j_k\in[f]$.
    We can define the type $\tau(j_1,\dots,j_k)\in\mathcal{N}^f$ as a string where the $i$-th value equals the number of $i$s in $(j_1,\dots,j_k)$ [for example, $\tau(1,2,1,1,3)=(3,1,1,0,\dots)$]. Then, it is immediate that 
    \begin{align}
        \left|\int_{T_\eps} dy\,q(y)\,y_{j_1}\cdots y_{j_k}\right|\le \eps^{1+(k-2)\alpha}
    \end{align}
    and it is nonzero only if  $\tau(j_1,\dots,j_k)$ consists of even numbers only. In particular, $\left|\int_{T_\eps} dy\,q(y)\,y_{j_1}\cdots y_{j_k}\right|=0$ as long as $k=3$.
    Finally, substituting the above bounds and the assumption (\ref{eq:derivativebound}) into Eq.~(\ref{eq:progerr1}), we have:
    \begin{align}
        \eps_{\rm prog}&\le \epsilon\sum_{j=1}^f\frac12\left\|\frac{\partial^2\map{U}_x}{\partial x_{j}^2}\right\|_{\diamond}+\sum_{k=4}^{\infty}\sum_{j_1,\dots,j_k=1}^{f}\left|\int_{T_\eps} dy\,q(y)y_{j_1}\cdots y_{j_k}\right|(2h)^{k}/2+o(\eps)\\
        &\le \epsilon\cdot\frac12\left\|\frac{\partial^2\map{U}_x}{\partial x_{j}^2}\right\|_{\diamond}+\frac{\epsilon^{1-2\alpha}}{2}\cdot\sum_{k=4}^{\infty}(2fh\epsilon^\alpha)^k+o(\eps)\\
        &= \epsilon\cdot\frac12\left\|\frac{\partial^2\map{U}_x}{\partial x_{j}^2}\right\|_{\diamond}+\frac{\epsilon}{2}\cdot O(\epsilon^{2\alpha})+o(\eps)\\
        &= \epsilon\cdot\frac12\left\|\frac{\partial^2\map{U}_x}{\partial x_{j}^2}\right\|_{\diamond}+\frac{\epsilon^{1-2\alpha}}{2}\cdot\sum_{k=4}^{\infty}(2fh\epsilon^\alpha)^k+o(\eps)\\
        &= \epsilon\cdot\frac12\left\|\frac{\partial^2\map{U}_x}{\partial x_{j}^2}\right\|_{\diamond}+o(\eps).
    \end{align}
\end{proof}

The performance of the concrete construction is an immediate consequence of Lemma \ref{lem:estprog}:
\begin{prop}
    The programming protocol based on sine-state-estimation (cf.~Section \ref{subsec:programmingunitary}) has an error upper bounded by $(1/D^2)\cdot\frac {f\pi^2}2\left\|\frac{\partial^2\map{U}_x}{\partial x_{j}^2}\right\|_{\diamond}+o(1/D^2)$.
\end{prop}
\begin{proof}
    For the sine-state-estimation-based protocol, the conditional probability density function is  
    \begin{align}      p(\hat{x}|x)&=\prod_{j=1}^f\left|\braket{\phi_{\hat{x}_j}}{\sin_{x_j}}\right|\\        &=\prod_{j=1}^f\left|\sqrt{2/D^2}\sum_{m=0}^{D-1}\sin\frac{(m+1/2)\pi}{D}e^{imy_j}\right|^2\quad y_j:=\hat{x}_j-x_j.
    \end{align}
    We now verify that all conditions in Lemma \ref{lem:estprog} are satisfied, and applying the lemma yields the desired result.
    It is straightforward to see that $p(\hat{x}|x)$ can be expressed as $q(y)$ for $y=\hat{x}-x$ (condition 1) and $q(y)=q(-y)$ for small $y$ (condition 3). The estimation error is evaluated to be (cf. Ref.~\cite{buvzek1999optimal}) 
    \begin{align}
        \epsilon_{\rm est}=\frac{f\pi^2}{D^2}+o(D^{-2}).
    \end{align} 
    At last, the probability density function for our estimation protocol has a width of $1/D$. Therefore, for any $\alpha<1/2$, we have $1-\int_{T_\eps}q(y)=o(\epsilon_{\rm est})$ (condition 2). 
\end{proof}

\subsection{A coherent protocol}
\YY{This part is just for my own record. I am not sure if it fits this work.}

It is a bit disappointing that the best programming protocol is always a measure-and-operate one. Is there a coherent protocol that achieves the same (asymptotic) performance? Here we present one.

Consider a task to program a unitary rotation $e^{-i\rho t}$, where $t$ is some known, fixed parameter, and $\rho$ encodes $f$ free parameters. A coherent way of achieving this is to prepare $\mathcal{E}(\rho^{\otimes m})$ as the program state, where $\mathcal{E}$ is the optimal encoder for $m$-copy quantum states in \cite{yang2018compression}. $\mathcal{E}(\rho^{\otimes m})$ requires a hybrid memory of $(f/2)\log m$ (qu)bits in the leading order of $m$.

To execute the rotation on receiving the program, one first decodes by applying the optimal decoder $\mathcal{D}$ in \cite{yang2018compression}. This results in $\rho^{\otimes m}$ with an error vanishing in $m$. Then, one applies the quantum principal component analysis (qPCA) algorithm  \cite{lloyd2014quantum}. The qPCA algorithm realizes the transformation $\rho^{\otimes m}\to e^{-i\rho t}$ as desired, up to an error of $O(t^2/m)$. Therefore, the qPCA error requires $m\sim1/\eps$.

\subsection{Converse}
We can show a lower bound on the programming cost as
\begin{align}\label{eq:prog-converse}
        \log|M|\ge (f/2-\delta)\log(1/\epsilon)-O(1)
    \end{align} for arbitrary $\epsilon$-independent $\delta>0$. 
The proof of our lower bound on $\log|M|$, the dimension of the program, shall follow the steps below:
\begin{enumerate}
    \item Following the same route as in Ref.~\cite[Supplementary Information]{yang2020optimal}, we can show that $\log |M|\ge I_H(\map{U}_x^{\otimes m}(\Phi_m),d p_x)-4m\sqrt{2\epsilon}\log d_m-1$. Here $I_H(\map{U}_x^{\otimes m}(\Psi_m),d x)$ is the Holevo information of the ensemble generated by applying $\map{U}_x^{\otimes m}$, drawn by a distribution $dp_x$ over $\mathbb{T}^f$, to a suitable input state $\Phi_m$, and $d_m$ is the dimension of the (support of the) ensemble bounded as $d_m\le(m+d^2-1)^{d^2-1}$.
    \item We show that there exist choices of $\Phi_m$ and $dp_x$ such that $I_H(\map{U}_x^{\otimes m}(\Phi_m),dp_x)\ge  f \log m-O(1)$. 
    \item Combining the above steps, we get  $\log |M|\ge (f-4m\sqrt{2\epsilon}/(d^2-1))\log m-O(1)$. Choosing $m\sim \delta/\sqrt{\epsilon}$ yields the desired bound (\ref{eq:prog-converse}).
\end{enumerate}
Next, we describe how Step 2 can be achieved:
\begin{itemize}
    \item First, we choose $\Phi_m$ to be the ``YRC sine-shaped state" (i.e. the optimal program state in \cite{yang2020optimal}) and $dp_x$ to be a uniform distribution on a meshed subset $\mathbb{M}^f\subset\mathbb{T}^f$. The meshed subset $\mathbb{M}^f$ satisfy $|x-x'|\le C/m$ for any $x,x'\in \mathbb{M}^f$ and some constant $C>0$, and thus  $\log|\mathbb{M}^f|=f\log m+O(1)$. For this we need some reasonable assumption on $\set{P}_f$.
    \item Notice that $I_H(\map{U}_x^{\otimes m}(\Phi_m),d p_x)=I(X:Q)$, where $X$ is a classical register that stores $x$, and $Q$ is a quantum register on which the state $\map{U}_x^{\otimes m}(\Phi_m)$ lives. By data-processing inequality, we have $I_H(\map{U}_x^{\otimes m}(\Phi_m),d p_x)\ge I(X:X')$, where $\hat{X}$ is an estimate of $X$, obtained by applying a (quantum) estimator on $\map{U}_x^{\otimes m}(\Phi_m)$ that yields a member $X'$ of the mesh $\mathbb{M}^f$.
    \item In Ref.~\cite{yang2020optimal}, we have shown that $\Phi_m$ is the asymptotically optimal state for metrology (i.e., estimation of $\map{U}$ given $m$ uses). Applying the optimal covariant measurement yields an estimate $\hat{U}_x$ of $U_x$ with error (in diamond norm) $O(1/m^2)$. Since there is an one-to-one and onto function from $\mathcal{T}^f$ to $\mathbb{F}$, it yields an estimate $\hat{X}$ (which is now a continuous variable) with high precision (attaining the  Heisenberg limit of quantum metrology):
    \begin{align}\label{eq:error-est}
        P(|\hat{X}-X|>c\cdot 1/m)\le P_e(m)
    \end{align}
    for some $c\in(0,C/2)$ that depends on the geometry of $\mathbb{F}$ (in fact, one should choose $C$ per $c$) and some $P_e(m)$ that vanishes (faster than logarithmically) as $m\to\infty$.
    \item Next, we define a map $\hat{X}\to X'$ that maps the continuous-value estimate to the closest point in the mesh $\mathbb{M}^f$. By Eq.~(\ref{eq:error-est}), the error probability of the estimator $Q\to \hat{X}\to X'$ should therefore not be larger than $P_e(m)$. By Fano's inequality, we get:
    \begin{align}
        I(X:X')\ge H(X)-H_2(P_e(m))-P_e(m)(\log(|\mathbb{M}^f|-1)),
    \end{align}
    where $H(X)=\log|\mathbb{M}^f|$ and $H_2(p)=-p\log p -(1-p)\log(1-p)$ is the binary entropy. Since $\mathbb{M}^f$ is a mesh where the distance between the neighbouring points is $C/m$, we have $|\mathbb{M}^f|\ge (L/C\cdot m)^f$ for some dimension parameter $L>0$. Therefore, we have:
     \begin{align}
        I(X:\hat{X})\ge (f-P_e(m))\log(Lm/C)-o(1).
    \end{align}
    Summarising, we get 
    \begin{align}
        I_H(\map{U}_x^{\otimes m}(\Phi_m),d p_x)\ge f\log(m)-o(1)
    \end{align}
    for small $\delta>0$ as desired.
\end{itemize}


\section{Free entropy of Hermitian operators: Observable programming}
It's natural to expect that the free entropy of a Hermitian operator characterizes the optimal programming of the corresponding observables. We consider two different formulations of the observable programming problem.


\subsection{Free entropy of an observable}
The free entropy of a Hermitian observable has the same formula as the free entropy of a state.
\begin{equation}
    \chi(H) = \int_\mathbb{R}\int_\mathbb{R}\log|x-y|\dd\mu_H(x)\dd\mu_H(y)\ .
\end{equation}
We define the physical free entropy as,
\begin{equation}
    \chi_\mathrm{phy}(H;\epsilon):=\delta(H)\log\epsilon^{-1} + \chi_\mathrm{reg}(H)
\end{equation}
where $H$ is the free entropy dimension of $H$, and $\chi_\mathrm{reg}(H)$ is the finite part of $\chi(H)$. Explicitly, for a $d$-dimensional Hermitian observable,
\begin{equation}
    \chi_\mathrm{phy}(H;\epsilon) = \left(1-\sum_i \frac{g_i^2}{d^2}\right)\log\epsilon^{-1} + \frac{2}{d^2}\sum_{i < j}\log g_ig_j|E_i-E_j|
\end{equation}
where $i$ labels distinct eigenvalues and $g_i$'s denote the eigenvalue degeneracies.


\subsection{Programming spectral measurements}

The first formulation is to program the measurement that corresponds to the observable. Given the spectral decomposition $H=\sum_i E_i P_i$, we define the programming problem as programming the projective measurement set $P:=\{P_i\}_i$. We assume that the eigenvalues $\{E_i\}$ are provided classically.

This problem has been studied in the literature \cite{DAriano_2005,PrezGarca_2006}. We want a program and a universal POVM set $(\sigma_H,\Pi)$, which simulates any measurement encoded by the program. Let the program be $\sigma_H$ and the state to be measured be $\rho$. To simulate the born rule statistics, the most general procedure is
\begin{equation}
    \tr\, \mathcal{P}(\rho\otimes\sigma_H)\cdot\Lambda_i \approx \tr\,P_i\rho
\end{equation}
Rearranging could eliminate the need to explicitly model $\mathcal{P}$
\begin{equation}
    \tr\, \rho\otimes\sigma_H\cdot \Pi_i\approx \tr\,P_i\rho,\quad\Pi_i:=\mathcal{P}^\dagger(\Lambda_i)
\end{equation}
Since we demand the simulation to work all all $\rho$, we have
\begin{equation}
    \tr_H\, (\id\otimes\sigma_H\cdot \Pi_i)\approx  P_i,
\end{equation}
and we quantify the error using
\begin{equation}
    d(P,\Pi) := \max_\rho\sum_i|\tr\, \rho(P_i-\Pi_i)|,
\end{equation}
or other appropriate distance measures.

This problem was analyzed in \cite{DAriano_2005,PrezGarca_2006}, but the direct part is not matched by the converse, and we shall try to improve it. 

\subsubsection{Achievability}
Since any projective measurement is effectively a computational basis measurement up to a change of basis, we could reduce this problem to unitary programming. 

Firstly, we consider the case where $H$ is non-degenerate and all $\{P_i\}$'s are rank-$1$ projectors. Consider the program $\sigma_U$ that encodes the basis transformation $U$ from the basis $\{P_i\}$ to the computational basis $\ketbra{i}{i}$, such that 
$\mathcal{P}(\,\cdot\,\otimes\sigma_U)\approx U\,\cdot\,U^\dagger$ in diamond norm distance. It follows that
\begin{equation}
   \tr\,\rho P_i=\bra{i}U\rho U^\dagger\ket{i}\approx \tr\, (\mathcal{P}(\rho\otimes\sigma_U)\cdot\ketbra{i}{i})= \tr\,\rho\otimes\sigma_U\cdot\mathcal{P}^\dagger(\ketbra{i}{i}),\quad\forall\rho
\end{equation}
Therefore, any $\eps$-unitary programming scheme $(\sigma_U,\mathcal{P})$ can be turned into a rank-$1$ projective measurement programming scheme $(\sigma_H=\sigma_U,\{\mathcal{P}(\ketbra{i}{i})\}_i)$ with the same fidelity. 

It would be an overkill to demand $(\sigma_U,\mathcal{P})$ to be a universal unitary programming scheme for any $U\in U(d)$. It is sufficient to program only the unitaries that satisfy
\begin{equation}
     \tr\,\rho P_i=\bra{i}U\rho U^\dagger\ket{i},\quad\forall \rho
\end{equation}
that is
\begin{equation}
    P_i=U^\dagger\ketbra{i}{i}U\ .
\end{equation}
These are unitaries in the coset,
\begin{equation}
    U(d)\Big/(U(1)\times\cdots\times U(1))\ ,
\end{equation}
and its real dimension (d.o.f.) is $d^2-d$, which is less than $d^2$ for a general element of $U(d)$. This is because we don't care about the $d$ phases associated with each basis transformation. Using the unitary programming protocol, we could program such a $U$ to $\sigma_U$ s.t.
\begin{equation}
    |\sigma_U| = \frac{d^2-d}{2}\log\eps^{-1} + \O(1) = \frac{d^2}{2}\delta(H)\log\eps^{-1} + \O(1)\ . 
\end{equation}\\

Let us now consider the general case where $H=\sum_{a=1}^k E_a P_a$ could be degenerate and $\{P_a\}_{a=1}^k$, $k<d$, could contain projectors of any rank. We need a unitary $V^\dagger$ that maps from $P_i$ to some projector spanned by the computational basis elements,
\begin{equation}
     U^\dagger P_a U=\sum_{i\in I_a}\ketbra{i}{i}
\end{equation}
where $g_i$ is also the rank of $P_i$; and $\{I_a\}_{a=1}^k$ is a partition of the index set $[d]$, i.e., $\cup_{a=1}^k I_a=[d]$.

Since we don't care how $U$ acts within each subspace projected by $\sum_{i\in I_a}\ketbra{i}{i}$, the space of distinct $U$ is the coset,
\begin{equation}\label{eq:flag_manifold}
    U(d)\Big/(U(g_1)\times\cdots\times U(g_k))
\end{equation}
and its real dimension (d.o.f.) is 
\begin{equation}
    d^2-\sum_a g_a^2\, \stackrel{!}{=}\, d^2\delta(H)\ .
\end{equation}
Then the optimal unitary programming result implies that 
\begin{equation}
    |\sigma_H| = \frac{d^2}2\delta(H)\log\eps^{-1}+\O(1)\ .
\end{equation}
Hence, we could match $\chi_\mathrm{phy}(H;\eps)$ at the leading order controlled by the free entropy dimension, but it's unclear how to obtain the second order term.


\subsubsection{Converse}
We can adapt the converse for unitary programming to this context to show that the measurement programming size $\chi_\mathrm{phy}(H;\eps)$ achieved above is optimal at the leading order.

Suppose that $H$ is non-degenerate. We want to argue that any spectral measurement programming scheme can be turned into an appropriate unitary programming scheme. To do so, we reverse the argument for the direct part. Given $(\sigma_H,\Pi_i)$, such that
\begin{equation}
    \tr\,\rho P_i\approx\tr\,\rho\otimes\sigma_H\cdot\Pi_i
\end{equation}
We can construct a CP unital channel $\mathcal{P}^\dagger$ that prepares the POVM $\{\Pi\}_i$
\begin{equation}
    \mathcal{P}^\dagger(\,\cdot\,):=\sum_j \Pi_j\, \bra{j}\,\cdot\,\ket{j},\quad\mathrm{s.t.}\quad \mathcal{P}^\dagger(\ketbra{i}{i})=\Pi_i
\end{equation}
Its dual channel $\mathcal{P}$ is a CPTP map that satisfies,
\begin{equation}
    \tr\,\rho P_i\approx\tr\,\mathcal{P}(\rho\otimes\sigma_H)\cdot\ketbra{i}{i}
\end{equation}

Each set of $P=\{P_i\}_i$ is in 1-1 correspondence with a unitary transformation $U_P$ in
\begin{equation}
        U(d)\Big/(U(1)\times\cdots\times U(1))
\end{equation}
in the sense that
$P_i=U_P\ketbra{i}{i}U_P^\dagger$. Therefore, we have
\begin{equation}
    \tr\,U_P^\dagger\rho U_P\cdot\ketbra{i}{i}\approx\tr\,\mathcal{P}(\rho\otimes\sigma_H)\cdot\ketbra{i}{i}
\end{equation}
Let $\rho=P_j$,
\begin{equation}
    \tr\,\mathcal{P}(P_j\otimes\sigma_H)\cdot\ketbra{i}{i}\approx \delta_{ij}\implies \tr\,\mathcal{P}(P_j\otimes\sigma_H)\approx \ketbra{j}{j} = U^\dagger_P P_j U
\end{equation}
This implies that $(\sigma_H,\mathcal{P})$ forms a universal programming scheme for unitaries in
\begin{equation}
        U(d)\Big/(U(1)\times\cdots\times U(1))\ .
\end{equation}
Then the converse bound for unitary programming implies that 
\begin{equation}
    |\sigma_H|\ge \frac{d^2-d}{2}\log\eps^{-1} +\O(1) = \frac{d^2}{2}\delta(H)\log\eps^{-1}+\O(1)
\end{equation}

When the observable $H$ is degenerate, we follow the same argument as above. Starting from a measurement programming scheme $(\sigma_H,\Pi_a)$ yields
\begin{equation}
    \tr\,\mathcal{P}(P_a\otimes\sigma_H)\approx \sum_{i\in I_a}\ketbra{i}{i} = U^\dagger_P P_j U_P
\end{equation}
This implies that $(\sigma_H,\mathcal{P})$ forms a universal programming scheme for unitaries in
\begin{equation}
        U(d)\Big/(U(g_1)\times\cdots\times U(g_k))\ .
\end{equation}
Then the converse bound for unitary programming implies that 
\begin{equation}
    |\sigma_H|\ge \frac{d^2-\sum_a g_a^2}{2}\log\eps^{-1} +\O(1) = \frac{d^2}{2}\delta(H)\log\eps^{-1}+\O(1)
\end{equation}

\subsection{Programming observable expectation}
The second formulation of the problem goes one step further from simulating the measurement operators. We program the observable directly by asking for a universal processor that, upon receiving an observable program, can directly produce the empirical expectation value of the observable for any state given in copies. 

We ask for a program and a universal processor, $(\sigma_H,\mathcal{P})$, such that when given enough copies of any state $\rho$, it outputs a value that is $\eps$-close to the true expectation $\tr\, H\rho$:

For any $\rho$, there is a sufficiently large $n$ such that
\begin{equation}
   |\mathcal{P}(\sigma_H\otimes\rho^{\otimes n}) -\tr\, H\rho|\le\eps.
\end{equation}
Here we do not restrain the number of copies $n$ so that the optimal simulation error $\eps$ is only controlled by the size of the program.

It is unclear if the approach of programming the spectral measurements would be optimal for this formulation of the problem.


\subsubsection{Achievability}
We present a protocol that allows us to program an observable in a quantum memory of size equal to the free entropy $\chi_\mathrm{phy}(H;\sqrt{\eps})$. Since the free entropy of an observable is essentially the same as for a density operator, the idea is to turn $H$ into a density operator $\tilde H$ whose $n$-copies-compression can serve as a program for $H$.

We first shift the minimal eigenvalue of $H$ to zero, $H':=H-E_0\id$. Then we normalize it to make a density operator $\tilde H:=H'/\tr\, H'=(H-E_0\id)/(\tr\,H-E_0d)$. Let $n=\eps^{-2}$. Our program $\sigma_H$ for $H$ is given by the compression for $\tilde H^{\otimes n}$, which has a size of $\frac{d^2}{2}\chi_\mathrm{phy}(\tilde H;1/n)$. Simple calculations show that
\begin{equation}
    |\sigma_H|=\frac{d^2}{2}\chi_\mathrm{phy}(\tilde H;1/n) = d^2\chi_\mathrm{phy}(H;\eps) - \log (\tr\, H-E_0d)\ .
\end{equation}
Hence, the program size is characterized by the physical free entropy of $H$ up to a $H$-dependent constant.

Note that we miss a factor of $2$ here in comparison to the measurement programming size discussed previously. This is consistent because the error parameter is different. If one were to measure the simulated $P_i$ on multiple copies of $\rho$ to extract the statistics, one would incur the same error scaling $n\sim \eps^{-2}$ from the Chernoff bound. 

For simulation, we feed in $n$ copies of a target state $\rho^{\otimes n}$. Our processor first decompresses $\sigma_H$ and then it performs the swap test between $\tilde H$ and $\rho$ for $n$ rounds, to get an estimate of
\begin{equation}
    \tr\, \rho\tilde H = \frac{\tr\, \rho H-E_0d}{\tr\, H-E_0d}
\end{equation}
which in turns gives us an estimate for $\tr\,\rho H$. 

Note that the decompressed state $\mathcal{D}(\sigma_H)$ is only approximately $\tilde H^{\otimes n}$ with an error $n^{-\frac12}\sim\eps$. This leads to a systematic error of $\O(\eps)$. For $n$ rounds of the swap test, the statistical error is also $\O(n^{-\frac12})\sim\O(\eps)$. Hence, the total error of simulating $\tr\,\rho H$ is $\O(\eps)$ as demanded.

\subsubsection{Converse}
\JW{Any ideas for a no-programming theorem?}


%\bibliographystyle{apsrev4-2}

\bibliographystyle{unsrt}
\bibliography{free}

\end{document}







